\PassOptionsToPackage{xetex,dvipsnames}{xcolor}
\PassOptionsToPackage{xetex}{graphicx}
\documentclass[xetex]{beamer}  % ue6aue6a
\title[Wordnets for everyone]{Building wordnets for languages with few resources}
%\\[2ex] \bl{Picking Brains, Crossing Ideas}}
\author[Bond and many more]{Francis \textbf{Bond} 
  \\ and many, many more}
\institute[Palacký]{Department of Asian Studies,
\\  Palacký University, Olomouc, Czechia
\\[1.5ex] \url{<bond@ieee.org>}}
%\date{\href{http://2016.fossasia.org/}{QTLEAP} 2016-09-10} %%% FIXME
%\date{CLUNL 2017-12-18} %%% FIXME
\date[2024-10-21]{Teanga Seminar Series} %%% FIXME
\newcommand{\ili}[1]{\href{http://www.globalwordnet.org/ili/#1}{\url{#1}}}

%%% make it smaller
 % gs -sDEVICE=pdfwrite -dCompatibilityLevel=1.4 -dPDFSETTINGS=/printer -dNOPAUSE -dQUIET -dBATCH -sOutputFile=omw-teanga-small.pdf omw-teanga.pdf 



%\usepackage[hidelinks]{hyperref}
\hypersetup{ 
     colorlinks, 
     linkcolor={blue!50!black},
     citecolor={red!50!black},
     urlcolor={blue!80!black}
}


\usepackage{array}
\newcolumntype{H}{>{\setbox0=\hbox\bgroup}c<{\egroup}@{}}

\usepackage[round]{natbib}
\usepackage{fontenc}
\usepackage{polyglossia}
\setmainlanguage{english}
\usepackage{xeCJK}
%\setCJKmainfont{WenQuanYi Micro Hei}
% \setmainfont[Numbers=OldStyle]{TeX Gyre Pagella}
%\setCJKmainfont{Droid Sans Fallback}
\setCJKmainfont{Noto Sans CJK SC}
\usefonttheme{professionalfonts}
\setsansfont{Linux Libertine O}
\logo{NTU}
\logo{
\includegraphics[width=15mm]{pics/UP_logo_FF_horizont_en.png}}

 \usepackage{qtree,avm,expex}
\definelingstyle{cmn}{everygla=,
everyglb=\itshape,everyglc=\footnotesize}

\usepackage{relsize,xspace}


\usetheme{Boadilla}
\usecolortheme{seahorse}
\setbeamertemplate{navigation symbols}{}

\newcommand{\zhong}{\ensuremath{[\Big|]}}
% NUS Meeting Room 1 (MR 1) 	COM1-03-19
\newcommand{\bl}[1]{\textcolor{blue}{#1}}
%\newcommand{\sh}[1]{\textcolor{blue}{#1}}
\newcommand{\txx}[1]{\textbf{\textcolor{blue}{#1}}}
\usepackage{url}

%\newcommand{\wn}{\url}
\usepackage[e,f,j]{mtg2e}
\newcommand{\cmn}{\mtciteform}
\newcommand{\kor}{\mtciteform}
\renewcommand{\mtcitestyle}[1]{\textcolor{teal}{\textit{#1}}}
\usepackage{xcolor}
 \newcommand{\bs}{\textbackslash}   % backslash
 \newcommand{\cmd}[1]{{\bf \color{red}#1}}   % highlights command
 \newcommand{\wemp}[1]{\ul{#1}}   % highlight
 \newcommand{\emp}[1]{\textbf{\color{blue}{#1}}}   % highlight
 \newcommand{\rel}[1]{\textsc{\color{blue}{#1}}}
 \newcommand{\lnk}[1]{\rel{#1}}

 \newcommand{\lex}[1]{\textbf{\color{teal}{\textit{#1}}}}
 \newcommand{\lng}[1]{\textcolor{teal}{\textit{#1}}}
\newcommand{\wn}[3]{\lex{#1}\ensuremath{_{#2:#3}}}
\usepackage{relsize}
\newcommand{\wnida}[1]{\asort{\wnid{#1}}}
\newcommand{\wnid}[1]{\textsf{\smaller #1}}

\usepackage{mygb4e,url}
\usepackage[normalem]{ulem}
\newcommand{\ul}{\uline}
\newcommand{\ull}{\uuline}
\newcommand{\udl}{\dashuline}
\newcommand{\uwl}{\uwave}

\newcommand{\MyLogo}[1]{(#1)}

\newcommand{\sh}[1]{https://fcbond.github.io/sh-canon/#1.html}
\newcommand{\SCAN}{\href{\sh{scan}}{\textit{A Scandal in Bohemia}}\xspace}
\newcommand{\REDH}{\href{\sh{redh}}{\textit{The Redheaded League}}\xspace}
\newcommand{\HOUN}{\href{\sh{houn}}{\textit{The Hound of the
      Baskervilles}}\xspace}
\newcommand{\SPEC}{\href{\sh{spec}}{\textit{The Adventure of the Speckled Band}}\xspace}
\newcommand{\DANC}{\href{\sh{danc}}{\textit{The Adventure of the Dancing Men}}\xspace}
\newcommand{\FINA}{\href{\sh{fina}}{\textit{The Adventure of the Final Problem}}\xspace}
\newcommand{\NAVA}{\href{\sh{nava}}{\textit{The Adventure of the Naval Treaty}}\xspace}

\newcommand{\swn}{SentiWordNet\xspace}
\newcommand{\plwn}{plWordNet\xspace}
\newcommand{\plemo}{plWordNet~emo\xspace}
\newcommand{\pwn}{Princeton WordNet\xspace} 
\newcommand{\mlsc}{ML-SentiCon\xspace}
\newcommand{\wnop}{\textsc{Micro-WNOp}\xspace}
\newcommand{\ntumc}{NTU-MC\xspace}

% typesetting: metaphors
\newcommand{\defn}[1]{{\textbf{#1}}}
\newcommand{\metaphor}[1]{\textsc{#1}}
\newcommand{\word}[1]{\textit{#1}}
\newcommand{\concept}[1]{\textsc{#1}}
\newcommand{\letter}[1]{\texttt{#1}}
\newcommand{\sense}[2]{\word{#1}\textsubscript{$#2$}}
\newcommand{\synonym}[1]{\textit{#1}}


\newcommand{\tabspace}{\addlinespace[0.6em]}
\newcommand{\slipspace}{\addlinespace[0.1em]}
\newcommand{\tighteq}{\,{=}\,}

\usepackage{adjustbox}
\usepackage{array}
\usepackage{tikz}

\newcommand{\literalcolour}{Brown}
\newcommand{\relatedcolour}{Bittersweet}
\newcommand{\metaphoricalcolour}{Purple}
\newcommand{\conduitcolour}{Gray}

\newcommand{\literallabel}{Prototype}
\newcommand{\relatedlabel}{Associated}
\newcommand{\metaphoricallabel}{Metaphorical}
\newcommand{\relatededgelabel}{Metonymy}
\newcommand{\relatededgelabelshort}{Meton.}
\newcommand{\metaphoricaledgelabel}{Metaphor}
\newcommand{\conduitlabel}{{ }+{ }Conduit}
\newcommand{\conduitmetlabel}{\hphantom{\metaphoricallabel}\conduitlabel}
\newcommand{\conduitasslabel}{\hphantom{\relatedlabel}\conduitlabel}

\newcolumntype{R}[2]{%
    >{\adjustbox{angle=#1,lap=\width-(#2)}\bgroup}%
    l%
    <{\egroup}%
}
\newcommand*\rot{\multicolumn{1}{R{30}{1.5em}}}

\newcommand{\literal}{\makecell{\textcolor{\literalcolour}{\literallabel}}}
\newcommand{\related}[2]{\makecell{\textcolor{\relatedcolour}{\relatedlabel} \\ \textcolor{\relatedcolour}{(from \sense{#1}{#2})}}}
\newcommand{\onlymetaphorical}{\makecell{\textcolor{\metaphoricalcolour}{\metaphoricallabel}}}
\newcommand{\metaphorical}[2]{\makecell{\textcolor{\metaphoricalcolour}{\metaphoricallabel} \\ \textcolor{\metaphoricalcolour}{(from \sense{#1}{#2})}}}
\newcommand{\feature}[1]{\scalebox{.8}[1.0]{\texttt{#1}}}
\newcommand{\yesfeature}[1]{\textcolor{Green}{\feature{#1}}}
\usepackage[normalem]{ulem}
\newcommand{\nofeature}[1]{\textcolor{Red}{\sout{\feature{#1}}}}
\newcommand{\maybefeature}[2]{\textcolor{Green}{\feature{#1}} \textcolor{Orange}{\feature{#2}}}
\newcommand{\allmaybefeature}[1]{\textcolor{Orange}{\feature{#1}}}
\newcommand{\defaultdepth}{0pt}
\tikzset{definition/.style={align=left, shape=rectangle,draw=black, font={\footnotesize}, fill=black!0!white}}
\tikzset{invis/.style={align=left, shape=rectangle,draw=none, font={\footnotesize}, fill=none,text opacity=0}}
\tikzset{lit_definition/.style={anchor=center, align=left, shape=rectangle,draw=\literalcolour, font={\footnotesize}, label={[label distance=-.05cm]:{\footnotesize{\textcolor{\literalcolour}{\literallabel}}}}, fill=black!0!white}}
\tikzset{met_definition/.style={align=left, shape=rectangle,draw=\metaphoricalcolour, font={\footnotesize}, fill=black!0!white}}
\tikzset{ass_definition/.style={align=left, shape=rectangle,draw=\relatedcolour, font={\footnotesize}, fill=black!0!white}}
\tikzset{ghost/.style={ass_definition, dashed}}
\tikzset{ghost_core/.style={lit_definition, dashed}}

% , on layer=bg
\tikzset{split/.style={-, dashed, line width=.3mm, Gray}}
\tikzset{start/.style={->, line width=.3mm, \literalcolour}}
\tikzset{analogy/.style={->, line width=.3mm, \metaphoricalcolour}}
\tikzset{related/.style={->, line width=.3mm, \relatedcolour}}
\tikzset{unknown/.style={->, line width=.3mm}}

\tikzset{derivation/.style={node distance = .7cm and %1.9cm}}
1cm}}
\tikzset{tight/.style={node distance = .55cm and .5cm}}
\tikzset{close/.style={node distance = .6cm and %1.9cm}}
.3cm}}

\newenvironment{tightcenter}{%
  \setlength\topsep{0pt}
  \setlength\parskip{0pt}
  \begin{center}
}{%
  \end{center}
}

\newcommand{\featuregoesto}{\textcolor{\metaphoricalcolour}{$\rightarrow$}}
% space hacks
\setlength{\belowcaptionskip}{-5pt}
\everypar{\looseness=-1}

\usepackage{tikz}
\usepackage{tikz-dependency}
\newcommand{\horn}[1]{\ensuremath{\langle}\lex{#1}\ensuremath{\rangle}}

\usepackage{tikz-qtree}


\begin{document}
\maketitle
\lingset{lingstyle=cmn}
\begin{frame}{Outline}
   \tableofcontents
 \end{frame}


% % \begin{frame}{Self Introduction}

% % \begin{itemize}\addtolength{\itemsep}{-0.5ex}
% % %\item Francis Bond 
% % \item BA in Japanese and Mathematics 
% %   \hfil\raisebox{-4.5ex}[0mm][0mm]{\includegraphics[width=5em]{pwm.pdf}}
% % \item BEng in Power and Control %(Electrical Systems Engineering)

% % \item PhD in English on  \textit{Determiners and Number in English
% %   contrasted with Japanese,  as exemplified in Machine
% %   Translation} 
% % \item 1991-2006 NTT (Nippon Telegraph and Telephone)
% %   \begin{itemize}
% %   \item Japanese - English/Malay Machine Translation
% %   \item Japanese corpus, HPSG grammar and ontology (Hinoki)
% %   \end{itemize}
% % \item 2006-2009 NICT (National Inst. for Info. and Comm. Technology)
% %   \begin{itemize}
% %   \item Japanese - English/Chinese Machine Translation
% %   \item Japanese WordNet
% %   \end{itemize}
% % \item 2009- NTU (Nanyang Technological University)
% %   \begin{itemize}
% %   \item Abui, Chinese, Malay, Multilingual Wordnets
% %   \item HPSGs for Chinese, Indonesian, \ldots
% %   \item Multilingual Meaning Banks (Treebank + Sensebank)
% %   \end{itemize}

% % \end{itemize}
% % \end{frame}
  
% \section{Wordnet}
% \begin{frame}{WordNet}

% %\MyLogo{\citet{Miller:1998:foreword,_Fellbaum:1998}}

% \begin{itemize}
% \item WordNet is an open-source electronic lexical database of English,
%   originally developed at Princeton University
%   \\   \url{http://wordnet.princeton.edu/}
% \item Made up of four separate (but interlinked) semantic nets, 
%   for nouns, verbs, adjectives and adverbs
% \item WordNets exist for many languages, I have worked on:
%   \begin{itemize}
%   \item Japanese 
%   \item Bahasa Malay/Indonesian
%   \item Chinese (Mandarin and Cantonese)
%   \item Abui (a Papuan language of Alor)
%   \item Odenet (German)
%   \item Kristang (Portuguese/Malay Creole)
%   \item The shared \txx{Open Multi-lingual Wordnet} (150+ languages) 
%     \\ \url{http://omwn.org}
% \end{itemize}
% \end{itemize}
% \end{frame}

% \begin{frame}{Wordnet Structure}
% \begin{itemize}
%  \item Lexical items are categorised into $\sim$120K (v 2023) glossed \textbf{synsets} (= synonym sets $\approx$ concepts)
%    \\ wordnet has grown over time, and continues to do so
%  \item Lexical relations are at either the synset level or sense (=
%    combination of lexical item and synset) level 
%  \item Strongly lexicalist (originally):
%    \begin{itemize}
%    \item synsets only where words exist
%    \item but many multiword expressions ($\approx 50\%$)
%   \item (near) absence of frame semantics
%    \end{itemize}
%  \end{itemize}
% \end{frame}


% \begin{frame}{The synset for  \eng[tiger]{虎}}
% \includegraphics[width=0.9\textwidth]{wn-jpn-tiger.pdf}

% Here we show English and Japanese.
% \end{frame}

% \begin{frame}{Another view of \eng{cat} as a graph}
% \includegraphics[width=0.9\textwidth]{graph.png}

% From \url{https://github.com/aliiae/lexical-graph}
% \end{frame}

% % \begin{frame}{Psycholinguistic Foundations of WordNet}

% % %\MyLogo{}

% % \begin{itemize}
% % \item Strong foundation on hypo/hypernymy (lexical inheritance) based on
% %   \begin{itemize}
% %   \item   response times to sentences such as:
% %     \\      \eng{a canary \{can sing/fly,has skin\}}
% %     \\    \eng{a bird \{can sing/fly,has skin\}}
% %     \\  \eng{an animal \{can sing/fly,has skin\}}
% %   \item analysis of anaphora:
% %     \\  \eng{I gave Kim a novel but the \{book,?product,...\}
% %       bored her}
% %     \\  \eng{Kim got a new car. It has shiny \{wheels,?wheel nuts,...\}}
% %   \item selectional restrictions
% %   \end{itemize}
% % \item Is now often used to calculate \txx{semantic similarity}
% %   \begin{itemize}
% %   \item The shorter the path between two synsets the more similar they are
% %   \item Or the shorter the path to the nearest shared hypernym, \ldots
% %   \item these paths can be weighted by information structure
% %   \end{itemize}
% % \end{itemize}
% % \end{frame}

% % \begin{frame}{Word Meaning as a Graph}
% % \bigskip
% % \bigskip
% % \bigskip
% % \includegraphics[height=0.75\textheight]{Semantic_Net.eps}
% %                                                         golf_equipment   golf   
% %                                                         club
% % operator                         utility program        wood
% % driver_1                         driver_2               driver_3
% % designated_driver tailgater

% % nondriver  computer_science

% % \begin{itemize}
% % \item You need a very big graph to capture all meanings
% % \end{itemize}

% % \end{frame}

% \begin{frame}{Relations}
%   \begin{itemize}
%   \item Constitutive
%     \\  \wn{driver}{n}{1} \rel{hyponym} \wn{operator}{n}{1}
%     \\  \wn{driver}{n}{1} \rel{antonym} \wn{nondriver}{n}{1}
%     \\   \wn{driver}{n}{3} \rel{synonym} \wn{number one wood}{n}{1}	  
%    \\   \wn{golf club}{n}{1} \rel{meronym} \wn{clubhead}{n}{1}	  
%   \item Derivational
%     \\  \wn{driver}{n}{1} \rel{derived} \wn{drive}{v}{1} 
%   \item Descriptive
%     \\ \wn{driver}{n}{2} \rel{domain} \wn{computer science}{n}{1}
%   \item Implied (from e.g. definitions)
%     \\ \wn{driver}{n}{1} ``the \wn{operator}{n}{1} of a \wn{motor
%       vehicle}{n}{1}''
%     \\ \wn{driver}{n}{2} ``a \wn{program}{n}{3} that
%     \wn{determines}{v}{2} how a \wn{computer}{n}{2} will
%     \wn{communicate}{n}{1} with a \wn{peripheral device}{n}{1}''
%   \end{itemize}

% \end{frame}

% \begin{frame}{Usage}
%   \begin{itemize}
%   \item Wordnets are used for 
%     \begin{itemize}
%     \item studying word similarity 
%     \item describe the meaning of texts
%     \item test WSD algorithms
%     \item improve parse ranking
%     \item machine translation
%     \item puzzle generation and solving
%     \item joke generation
%     \item distinguishing senses for word embeddings
%     \item language learning
%     \end{itemize}
%   \item They are still useful in the age of LLMs
%     \begin{itemize}
%     \item As a check on quality
%     \item For explicit reasoning
%     \item For languages with insufficient language data
%     \item For teaching
%     \end{itemize}
%   \end{itemize}

% \end{frame}

% \subsection{Wordnets in the World}

% \begin{frame}{Wordnets in the world 2008-06}
%  Many wordnets, but few free: we decided to build an open wordnet for Japanese and needed other wordnets for cross-lingual disambiguation.
%  \begin{center}
%    \noindent\includegraphics[width=1.0\textwidth]{wns-2008.png}
%  \end{center}
% Green is free; Blue is research only; Brown costs money

% \end{frame}


% \begin{frame}{Wordnets in the world 2008-06}

%   \begin{itemize}
%   \item When we decided to build the Japanese wordnet, we wanted to
%     leverage existing work and use multiple languages to disambiguate
%   \item Many wordnets existed, but few were free and almost all came
%     in slightly (or radically) different formats.
%     \begin{itemize}
%     \item PWN had a loop!
%     \item GWG grid had bad encodings
%     \item There was \emp{no} wordnet that worked out of the box
%     \item  Different projects even used different names for pos
%       \\ e.g., Adjective ($a$ vs $j$); Adverb ($r$ vs $b$) 
%     \end{itemize}
%   \item There was no easy shared access
%   \item We had to massage the data for our project
%   \item  We wanted to try to save other people the trouble we
%     experienced, and make the normalized data available (where legally possible)
%   \end{itemize}
% \end{frame}




% \begin{frame}{Wordnets in the world 2011-06}
% Free wordnets for French, Catalan, Polish, Thai and \txx{Japanese}
% \begin{center}
%   \noindent\includegraphics[width=1.0\textwidth]{wns-2011.png}
% \end{center}
% Green is free; Blue is research only; Brown costs money
% \end{frame}


% \begin{frame}{Wordnets in the world 2012-06}
% Spanish freed (and Galician and Basque), Farsi, Norwegian, Swedish, \txx{Bahasa} (Malay and Indonesian)
% \begin{center}
%   \noindent\includegraphics[width=1.0\textwidth]{wns-2012.png}
% \end{center}
% Green is free; Blue is research only; Brown costs money
% \end{frame}

% \begin{frame}{Wordnets in the world 2013-06}
% \txx{Chinese} Open Wordnet (and Chinese wordnet)
% \begin{center}
%   \noindent\includegraphics[width=1.0\textwidth]{wns-2013.png}
% \end{center}
% Green is free; Blue is research only; Brown costs money
% \end{frame}

% \begin{frame}{The synset for  \eng[tiger]{虎}}
% \includegraphics[width=0.9\textwidth]{wn-jpn-tiger.pdf}

% % \begin{frame}{Wordnets in the world 2014-06}
% % Swedish
% % \begin{center}
% %   \noindent\includegraphics[width=1.0\textwidth]{WordNet-map-2014-06}
% % \end{center}
% % Green is free; Blue is research only; Brown costs money
% % \end{frame}



% % \begin{frame}{Wordnets in the world 2015-06}

% % \begin{center}
% %   \noindent\includegraphics[width=1.0\textwidth]{WordNet-map-2015-06}
% % \end{center}
% % Green is free; Blue is research only; Brown costs money
% % %\MyLogo{}
% % %\MyLogo{}
% % \end{frame}



% \begin{frame}{Wordnets in the world 2016-01}
% Swedish, Dutch, Icelandic, Lithuanian, Romanian, \ldots
% \begin{center}
%   \noindent\includegraphics[width=1.0\textwidth]{wns-2016.png}
% \end{center}
% Green is free; Blue is research only; Brown costs money
% %\MyLogo{}
% %\MyLogo{}
% \end{frame}



% \begin{frame}{}
% \frametitle{Why do we want so many?}
% \begin{itemize}
% \item Every new wordnet makes the network much richer $n(n-1)$ lexicons!
% \item Multilingual disambiguation makes it easier to add new languages
% \item New languages add new phenomena 
%   \begin{itemize}
%   \item and new synsets/concepts
%   \item and new relations
%   \end{itemize}
% \item More users mean more bugs found
% \item New approaches can be shared
%   \begin{itemize}
%   \item UKB  (does graph-based WSD) \hfill Basque
%   \item Wordnet glosses  (disambiguates wordnet definitions) \hfill Princeton 
%   \item Logical forms  (gives LF for wordnet definitions) \hfill USC/ISI
%   \item UKB + Wordnet glosses + Logical forms  (better WSD) \hfill Bulgaria
%   \item[\ldots]
%   \end{itemize}
% \end{itemize}
% \end{frame}

% \begin{frame}{How did we open the data?}
%   \begin{itemize}
%   \item Leading by example (we made open wordnets for Japanese, Bahasa, Chinese)
%   \item Appeal to self-interest: we showed that open resources are cited more
%      \citep{Bond:Paik:2012}
%   \item Simple format for sharing (tsv) --- I wrote many converters
%     \\ Many checks for ill-formed lexicons
%   \item Public praise for freed resources
%   \item Private persuasion for non-open resources
%   \item Open website 
%     \begin{itemize}
%     \item online interface (with statistics on coverage)
%     \item downloadable in multiple formats
%     \item linked to other resources (sentiment, time, SUMO, \ldots)
%     \end{itemize}
%   \item Used by other projects: Google Translate; 
%     Natural Language Toolkit (NLTK); Babelnet
%   \end{itemize}
% \end{frame}


% \begin{frame}{Wordnets in the world 2017-12}
% Added: Moroccan Arabic, Abui, Myanmar, Sesotho, Tswana, Venda, Xhosa, Zulu;
% \\Ancient: Greek, Latin, Sanskrit; Freed: German, Estonian, Romanian, Czech
% \begin{center}
%   \noindent\includegraphics[width=1.0\textwidth]{wns-2017.png}
% \end{center}
% %Green is free; Blue is research only; Brown costs money
% %\MyLogo{}
% %\MyLogo{}
% \end{frame}

% \begin{frame}{Wordnets in the world 2019-06}
%   Added: Cantonese, Coptic, Russian, Mongolian, Danish
%   \\ Expanded: English
% \begin{center}
%   \noindent\includegraphics[width=1.0\textwidth]{wns-2019.png}
% \end{center}
% \end{frame}

% \begin{frame}{Wordnets in the world 2021-01}
%   Added: Uzbek, Turkish;
% \\  Richer inventory of relations, better documentation
% \begin{center}
%   \noindent\includegraphics[width=1.0\textwidth]{wns-2021.png}
% \end{center}
% \end{frame}

% \begin{frame}{Wordnets in the world 2023-01}
%   Added: Ukrainian;
% \\  Richer inventory of relations, better documentation
% \begin{center}
%   \noindent\includegraphics[width=1.0\textwidth]{wns-2023.png}
% \end{center}
% \end{frame}



% \begin{frame}{Current state}
%   \begin{itemize}
%   \item 120,000+ synsets from the Open English Wordnet
%     \begin{itemize}
%     \item 2,000 added locally 
%     \end{itemize}
%   \item 35 curated wordnets with 2,000,000 senses
%   \item 250 languages with automatically created senses from
%     wiktionary (at least 100 senses per language)
%   \item 1,200 languages with senses mapped from various Swadesh lists
%     (around 100-200 senses/language)
%   \end{itemize}
% \end{frame}



% % \begin{frame}{Open Wordnets in the world 2017-01}


% %   \begin{large}
% %     Albanet; Arabic WordNet (AWN v2); BulTreeBank Wordnet (BTB-WN);
% %     Chinese Open Wordnet; Chinese Wordnet (Taiwan); DanNet; Greek
% %     Wordnet; Princeton WordNet; Persian Wordnet; FinnWordNet; WOLF
% %     (Wordnet Libre du Français); Hebrew Wordnet; Croatian Wordnet;
% %     IceWordNet; MultiWordNet; ItalWordnet; Japanese Wordnet;
% %     Multilingual Central Repository (Spanish, Galician, Catalan,
% %     Basque); Wordnet Bahasa (Malay, Indonesian); Open
% %     Dutch WordNet; Norwegian Wordnet (Bokmal); Norwegian Wordnet
% %     (Nynorsk); plWordNet; OpenWN-PT; Romanian Wordnet; Lithuanian
% %     WordNet; Slovak WordNet; sloWNet; Swedish (SALDO); Thai Wordnet
% %   \end{large}
%         \nocite{Ruci:2008,Black:Elkateb:Rodriguez:Alkhalifa:Vossen:Pease:Bertran:Fellbaum:2006,Borin:Forsberg:Lonngren:2013,Pedersen:2009,Simov:Osenova:2010,Gonzalez-Agirre:Laparra:Rigau:2012,Pociello:Agirre:Aldezabal:2011,Wang:Bond:2013,Huang:Hsieh:Hong:Chen:Su:Chen:Huang:2010,Pedersen:2009,_Fellbaum:1998,Stamou:Nenadic:Christodoulakis:2004,Sagot:Fiser:2008,Ordan:Wintner:2007,Noor:Sapuan:Bond:2011,Isahara:Bond:Uchimoto:Utiyama:Kanzaki:2008,Montazery:Faili:2010,Linden:Carlson:2010,Garabik:Pileckyte:2013,Vossen:Postma:2014,Piasecki:Szpakowicz:Broda:2009,Paiva:Rademaker:Melo:2012,Tufis:Ion:Bozianu:Ceausu:Stefanescu:2008,Fiser:Novak:Eejavec:2012,Borin:Forsberg:Lonngren:2013,Thoongsup:Charoenporn:Robkop:Sinthurahat:Mokarat:Sornlertlamvanich:Isahara:2009,Pianta:Bentivogli:Girardi:2002,Oliver:Sojat:Srebacic:2015,Raffaelli:2008,Toral:Bracal:Monachini:Sora:2010}

        
% % \medskip
% % And in the pipeline: Irish, Ancient Greek, Latin, Myanmar, Abui, \ldots
% %\end{frame}

% \begin{frame}{}
% \frametitle{The Open Multilingual Wordnet 1.0}
% \begin{itemize}
% \item Defined a minimal shared format (incrementally useful)
%   \\ based on the hierarchy of PWN~3.0 (most wordnets are translations)
%   \begin{itemize}
%   \item \txx{name, license, URL} (metadata)
%   \item \txx{synset \& lemma} pairs (linked to PWN)
%   \item later also \txx{definitions} and \txx{examples} 
%   \end{itemize}
% \item Wrote conversion scripts to extract this information
%   \begin{itemize}
%   \item Provided feedback when there were errors
%   \item Occasionally suggested improvements
%   \end{itemize}
% \item Encouraged people to choose an open license 
%   \\ we are deriving a   new resource and redistributing 
%   \\ --- this needs permission
% \item Linked each resource to its canonical citation
%   \\ Encouraged people to cite the papers when they used the wordnets
% \end{itemize}
% Mostly done by me on weekends, some help from Lars Nygaard, John McCrae and of course the individual projects.
% \end{frame}

% \begin{frame}{What else did we do?}
%   \begin{itemize}
%   \item Minimal evaluation
%     \begin{itemize}
%     \item Coverage of core concepts
%       \\ made the core concept list more easily available
%     \end{itemize}
%   \item Created a common search interface
%     \begin{itemize}
%     \item with links to other external resources: SUMO, sentiment, temporal, \ldots
%     \end{itemize}
%   \item Provided downloads in standard formats (LMF, LEMON)
%   \item Automatically created wordnets from Wiktionary and CLDR
%     \begin{itemize}
%     \item shared this data freely 
%     \end{itemize}
%   \item Created a Python API to access the data in NLTK \citep{Bird:Klein:Loper:2009}
%     \begin{itemize}
%     \item Extended the NLTK corpus class to add information about recommended citation and licenses (for all corpora, not just the wordnets)
%     \item New wordnets and updates are still being added
%     \end{itemize}
%   \end{itemize}

% \end{frame}



% % \subsection{NLTK}


% % \begin{frame}[fragile]{Python API: NLTK}

% % \begin{verbatim}
% % >> import nltk
% % >>> from nltk.corpus import wordnet as wn
% % >>> wn.synsets('dog')
% %  [Synset('dog.n.01'), Synset('frump.n.01'), 
% %   Synset('dog.n.03'), Synset('cad.n.01'), 
% %   Synset('frank.n.02'),  Synset('pawl.n.01'), 
% %   Synset('andiron.n.01'), Synset('chase.v.01')]

% % >>> wn.synsets('dog')[0].definition()
% %  'a member of the genus Canis (probably descended 
% %   from the common wolf) that has been domesticated 
% %   by man since prehistoric times; occurs in many breeds'
% % \end{verbatim}
% % \end{frame}

% % \begin{frame}[fragile]{Python API: NLTK}

% % \begin{verbatim}
% % >>> wn.synsets('dog')[0].lemmas(lang='jpn')
% %  [Lemma('dog.n.01.イヌ'), Lemma('dog.n.01.ドッグ'), 
% %   Lemma('dog.n.01.洋犬'), Lemma('dog.n.01.犬')]

% % >>> wn.synsets('dog')[0].lemmas(lang='zsm')
% %  [Lemma('dog.n.01.anjing')]

% % >>> wn.synsets('dog')[0].hypernym_paths()[0]
% %  [Synset('entity.n.01'), Synset('physical_entity.n.01'), 
% %   Synset('object.n.01'), Synset('whole.n.02'), 
% %   Synset('living_thing.n.01'), Synset('organism.n.01'), 
% %   Synset('animal.n.01'), Synset('chordate.n.01'), 
% %   Synset('vertebrate.n.01'), Synset('mammal.n.01'), 
% %   Synset('placental.n.01'), Synset('carnivore.n.01'), 
% %   Synset('canine.n.02'), Synset('dog.n.01')]
% % \end{verbatim}
% % \end{frame}




% % \subsection{Chinese Open Wordnet}
% % \begin{frame}{Chinese Open Wordnet (COW)}

% %   \begin{itemize}
% %   \item \href{http://compling.hss.ntu.edu.sg/cow/}{Chinese Open
% %       Wordnet --- 汉语开放词网
% % }
% %   \item 42,315 synsets
% %   \item 79,812 senses
% %   \item 61,536 unique words
% %   \item built by Shan Wang and Francis Bond
% %     \begin{itemize}
% %     \item with cooperation from Chinese Wordnet (South Eastern
% %       University), Bilingual Ontological Wordnet (Academica Sinica) and OMW
% %     \end{itemize}
% %   \item new release coming soon with pronouns, exclamatives,
% %     classifiers, chengyu and pinyin.
% %   \item We would like more collaborators, \ldots
% %   \end{itemize}

% % \end{frame}


% % \begin{frame}{Chengyu}
% %   \begin{itemize}
% %   \item include pinyin for entries, with multiple variants
% %     \begin{itemize}
% %     \item \cmn[woman]{女人}
% %     \item \cmn{nǚrén}
% %     \item \cmn{nü3ren2}
% %     \item \cmn{nv3ren2}
% %     \item \cmn{nüren}
% %     \item \cmn{nvren}
% %     \end{itemize}
% %   \item pinyin added from open source lexicons, with a confidence
% %     based on whether we found the whole word and if it was ambiguous
% %     or not
% %   \end{itemize}

% % \end{frame}

% % \begin{frame}{Interjections}
% %   \begin{avm}
% % \[ \wnida{\wnid{80000000-x} (pain)}
% % eng-lemmas & \parbox[t]{0.6\textwidth}{\eng{ah, oh, ouch, ow, wow, yipe, yow}} \\ 
% % cmn-lemmas & \parbox[t]{0.6\textwidth}{\cmn{哎呀, 啊呀, 哎哟, 啊哟, 哎哟喂, 疼}} \\ 
% % definition & \parbox[t]{0.6\textwidth}{an expression that is uttered to show physical hurt or pain} \\ 
% % exemplifies & \wnid{15167474-n} (utterance) \\ 
% % see also & \wnid{14322699-n} (pain, hurting) \]
% % \end{avm} \\
% % \begin{avm}
% % \[ \wnida{\wnid{80000002-x} (checkmate)}
% % eng-lemmas & \eng{checkmate, mate}  \\ 
% % cmn-lemmas & \eng{将死}  \\ 
% % definition & \parbox[t]{0.6\textwidth}{an expression that is uttered during a game of chess to declare that the final winning move has taken place} \\ 
% % exemplifies & \wnid{15167474-n} (utterance) \\ 
% % see also & \wnid{00167764-n} (checkmate) \]
% % \end{avm}
% % \end{frame}

% % \begin{frame}{Classifiers}
% %   \begin{small}
% % \begin{avm}
% % \[ \wnida{80000003-x }
% % cmn-lemmas & \textcolor{teal}{把} \cmn{bǎ} \\ 
% % \parbox[t]{0.2\textwidth}{definition} & \parbox[t]{0.6\textwidth}{a sortal classifier used with tools and objects with a handle, such as hammers, brooms, guitars or teapots} \\ 
% % exemplifies & \wnid{06308436-n} (classifier) \\
% % classifies & \wnid{03481172-n} (hammer) \\
% % classifies & \wnid{02906734-n} (broom) \\
% % classifies & \wnid{04398044-n} (teapot) \]
% % \end{avm} \\

% % % \begin{avm}
% % % \[ \wnida{80000004-x }
% % % lemmas & \textcolor{teal}{根} \jpn{gen} \\ 
% % % definition & \parbox[t]{16em}{a sortal classifier used for long slender objects, such as a banana, a pillar, a sausage or a needle} \\ 
% % % domain usage & \wnid{06308436-n} (classifier)  \]
% % % \end{avm} \\

% % \begin{avm}
% % \[ \wnida{80000004-x }
% % jpn-lemmas & \textcolor{teal}{振} \jpn{furi} \\ 
% % \parbox[t]{0.2\textwidth}{definition} & \parbox[t]{0.6\textwidth}{a sortal classifier used for weapons with a blade, such as knifes, swords or daggers} \\ 
% % exemplifies & \wnid{06308436-n} (classifier) \\
% % classifies & \wnid{03624134-n} (knife) \\
% % classifies & \wnid{04373894-n} (sword) \\
% % classifies & \wnid{03158885-n} (dagger) \]
% % \end{avm}
    
% %   \end{small}
% % \end{frame}

% % \begin{frame}{Idioms}
% %   \begin{avm}
% %     \[ \wnida{80000005-a}
% % cmn-lemmas & \textcolor{teal}{根深蒂固} \cmn{gēnshēndìgù} \\ 
% % \parbox[t]{0.2\textwidth}{definition$_{eng}$}
% % & \parbox[t]{0.6\textwidth}{deep-rooted (lit: roots deep stem strong)}
% % \\ 
% % \parbox[t]{0.2\textwidth}{definition$_{cmn}$} & \parbox[t]{0.6\textwidth}{基础深厚，难以动摇} \\ 
% % exemplifies & \wn{chengyu}{n}{1}  \\
% % similar-to  & \wn{deep-rooted}{a}{1} \\
% %  example & \[ essay$_{176}$ & 罪魁祸首自然是那些\wemp{根深蒂固}的\\
% %                    &  错误和持续的恶性循环。 \\
% %                     news$_{101303}$ & 另外，对再次被日本统治的警戒 \\
% %                    &   论也\wemp{根深蒂固}。 \\
% %                     news$_{101421}$ & 对贝鲁斯柯尼的搜查是成为促使 \\
% %                    &  这个战后政治崩溃之契机的净化 \\
% %                    &  作战的一环，问题\wemp{根深蒂固}。\] \\
% %        xref & \[ Chengyu500 & \#37 p122  \]
% % \]
% %   \end{avm}
% % \end{frame}
% \subsection{The Japanese Wordnet}

% \begin{frame}{The Japanese Wordnet}
% \begin{itemize}
% \item Initially assume that the semantic structure is the same
%   \begin{itemize}
%   \item \lex{dog} $\subset$ \eng{animal}
%   \item[$\Rightarrow$] \lex{犬} $\subset$ \lex{動物}
%   \end{itemize}
% \item Added Japanese words to Princeton wordnet synsets
% \end{itemize}
% \begin{tabular}{ccrrrl}
%   Date       & Ver & Concepts & Words & Senses  & Misc \\ \hline
%   2009-02-28 & 0.90 &  49,190 & 75,966 & 156,684 & initial release \\
%   2009-08-31 & 0.91 &  50,739 & 88,146 & 151,831 & linked to SUMO \\
%   2009-11-16 & 0.91 &  49,655 & 87,133 & 146,811  & \\
%   2010-03-05 & 1.00 &  56,741 & 92,241 & 157,398 & + definitions, examples \\
%   2010-10-22 & 1.10 &  57,238 & 93,834 & 158,058 & \\
%   2012-01-06 &      &         &        &          & Japanese Semcor \\
%   2014-02-06 &      &         &        &          & NLTK module \\
%   \textbf{2022-05-22} & 2.0  &  58,039 & 89,902 & 147,207 & 386,222  variants \\
% \end{tabular}
% \end{frame}

% \begin{frame}{Japanese Wordnet 2.0}
%   \begin{itemize}
%   \item This release uses a new format (GWA XML 1.0)
%     \begin{itemize}
%     \item Allows for variant forms
%     \item Structure can be different from PWN
%       \begin{itemize}
%       \item 788 new concepts (synsets)
%       \end{itemize}
%     \item Includes frequencies from JSemCor and NTU-MC
%     \end{itemize}
%   \item A long slog by Takyuki Kuribayashi to add the forms
    
%   \end{itemize}
% \end{frame}


% \begin{frame}{Orthographic Variants}

%   \begin{itemize}
%   \item synsetID=14728724-n (Eng: protein)
%     \\ プロテイン, 蛋白質, タンパク, たんぱく質, 蛋白, タンパク質
%     \\ \hspace{20em} ↓
%     \\ \textbf{蛋白質} (タンパクシツ, たんぱく質, タンパク質, たんぱくしつ)
%     \\ \textbf{蛋白} (タンパク, たんぱく)
%     \\ \textbf{プロティン} (プロテイン, ぷろてぃん)
%   \item synsetID=02765464-v (Eng: absorb, take in)
%     \\ 呑みこむ, 呑込む, 吸引, 吸い込む, 吸収
%     \\\hspace{20em} ↓
%     \\ \textbf{吸い込む} (スイコム, 吸込む, 吸いこむ, すいこむ)
%     \\ \textbf{吸収} (キュウシュウ, きゅうしゅう)
%     \\ \textbf{吸引} (キュウイン, きゅういん)
%     \\ \textbf{飲み込む} (ノミコム, 飲込む, 呑み込む, 呑込む, 呑みこむ, のみ込む, のみこむ)
%   \end{itemize}
% \end{frame}

% \begin{frame}{Other extensions}
%   \begin{itemize}
%   \item Added pronouns and demonstratives
%     \\ \lex{私, この, その, あの, どの}
%   \item Added classifiers
%     \\ \lex{人, 台, 匹, 回}
%   \item Added 4-character idioms
%     \\ \lex{五十歩百歩}
%   \item Added corpus-based examples and sense frequency
%     \\ \lex{対策}$_3$, \lex{策}$_3$, \lex{措置}$_2$,\lex{方略, 方策, 術, 打つ手}
%     ``step, measure''
%    \item Added exclamatives
%     \\ \lex{王手, ヨイショ, お早うございます}
%   \end{itemize}
% \end{frame}


% \begin{frame}{Ontological Differences: Temperature}

%   \begin{itemize}
%   \item English uses the same words for temperature experienced by
%     touching or as a general feeling:
% \\ \horn{cold, cool, warm, hot}
% \item Japanese distinguishes
%   \begin{itemize}
%   \item the feeling: \horn{寒い, 涼しい, 温かい, 暑い}
%   \item to-touch:  \horn{\scalebox{2}[1]{冷い}, 暖か, 熱い}
%   \end{itemize}
% \item English uses a single word for water of any temperature:
%   \lex{water}
% \item Japanese uses different words for cold (non-hot) water and hot
%   water: \lex{水} vs \lex{湯}
% \item We are doing our best to add native Japanese concepts, even if not lexicalized in English
% \end{itemize}
% \end{frame}

% \begin{frame}{Ontological Differences: Kin}
%   \begin{tikzpicture}
% \tikzset{level distance=12ex}
% \Tree [.\node(ja){\begin{avm}% 
%     \[ \wnida{jwn:3478-n}
%     lemma & \eng{姉妹兄弟, 兄弟} \\ 
%     \]\end{avm}};
%       [.\node(jb){\begin{avm}% 
%     \[ \wnida{jwn:6597-n}
%     lemma & \eng{兄弟, 男兄弟} \\ 
%     \]\end{avm}}; 
%  [.\node(je){\begin{avm}% 
%     \[ \wnida{jwn:6897-n}
%     lemma & \eng{兄} \\ 
%     \]\end{avm}}; ]
%  [.\node(jf){\begin{avm}% 
%     \[ \wnida{jwn:5957-n}
%     lemma & \eng{弟} \\ 
%     \]\end{avm}}; ]
% ]
%       [.\node(jc){\begin{avm}% 
%     \[ \wnida{jwn:8745-n}
%     lemma & \eng{姉妹} \\ 
%     \]\end{avm}};  [.\node(jg){\begin{avm}% 
%     \[ \wnida{jwn:2897-n}
%     lemma & \eng{姉} \\ 
%     \]\end{avm}}; ]
%  [.\node(jh){\begin{avm}% 
%     \[ \wnida{jwn:2957-n}
%     lemma & \eng{妹} \\ 
%     \]\end{avm}}; ]]
%          ]
%     \end{tikzpicture}

%     \begin{itemize}
%     \item More accurately models Japanese
%     \item We aim to cover at least the TUFS basic vocabulary
%     \end{itemize}


  
% \end{frame}

% \subsection{Korean Wordnets}
% \begin{frame}{Korean Wordnets}

%   \begin{itemize}
%   \item There is a good wordnet built at Pusan National University
%     \citep{Yoon:Hwang:Lee:Kwon:2009}, but unfortunately it is not open
%   \item There is a small wordnet built from the Tokyo University of
%     Foreign Studies Basic Vocabulary list
%     \begin{itemize}
%     \item It is free (CC By 4.0)
%     \item 527 concepts; 1,011 words; 1,698 example sentences
%       \begin{avm}
%         \[ \wnida{10694258-n}
%           en-lemmas & \cmn{teacher} \\
%           ko-lemmas & \textcolor{teal}{선생님} \cmn{seonsaengnim} \\ 
%           \parbox[t]{0.2\textwidth}{definition} & \parbox[t]{0.6\textwidth}{a person whose occupation is teaching} \\
%           ko-example & \ul{선생님}, 이 보고서 말인데요, 마감이 내일까지죠?
%           \]
%         \end{avm} 
%       \item Not really big enough to do much with
%       \item and biased to teaching --- should the entry be just 선생?
%     \end{itemize}
    
%   \end{itemize}
% \end{frame}

% \subsection{The Open English Wordnet}
%   \begin{frame}{The Open English Wordnet}
%       \begin{itemize}
%       \item The successor to the Princeton Wordnet (which is no
%         longer developed)
%       \item New vocabulary (thousands of new words and senses)
%       \item New relationships (\lnk{taboo}, \lnk{agent}, \lnk{female-form})
%       \item Many, many bug fixes
%       \item Pronunciation (in IPA)
%       \end{itemize}

%   \end{frame}

  
% % \begin{frame}{Open Multilingual Wordnet 2.0}

% %   \begin{itemize}
% %   \item Instead of a single hierarchy (PWN~3.0) it is the union of
% %     multiple hierarchies
% %     \begin{itemize}
% %     \item Allow links (relations) and nodes (synsets) not in PWN
% %       \begin{itemize}
% %       \item Merge nodes that link to the same Interlingual Index
% %         record
% %       \item Map all semantic relations to a known subset
% %       \item Document them
% %       \item Add more as the need is shown
% %       \end{itemize}
% %       \item Now we can better add in differently structured wordnets
% %         like Polish, Dutch or Gaelic
% %     \end{itemize}
% %   \item The goal is to have a useful subset, without covering
% %     everything from every wordnet (too much work)
% %   \item The conversion (once again) reveals issues in individual wordnets
% %   \item We hope to again link to NLTK, this time with more
% %     applications (e.g., WSD) and links to corpora 
% %   \end{itemize}
% % \end{frame}


% % \subsection{NTU-MC}
% % \begin{frame}{NTU multilingual corpus}

% % We use it to describe the meaning of texts

% % \begin{itemize}
% % \item Small, deeply analysed corpus
% % \begin{itemize}
% % \item 6,000+ sentences x 3 languages (cmn, eng, jpn); 2,000 in Indonesian
% %     \begin{itemize}
% %     \item Mainichi Newspaper (NICT translations)
% %     \item Sherlock Holmes (\sh{DANC}, \sh{SPEC}, \sh{REDH}[eng]) 
% %     \item Cathedral and the Bazaar (many languages)
% %     \item Spider's thread 「蜘蛛の糸」
% %     \item Singapore Tourist data (plus Korean,  Viet, Indo)
% %     \end{itemize}
% %   \item \bl{Hand alignment, WordNet tagging}, Treebanking, 
% %     Sentiment, \ldots
% %   \end{itemize}
% % \item Revised the DB structure; second round of wordnet tagging
% %   \begin{itemize}
% %   \item Slightly behind schedule
% %   \item Wordnet extension is being done in parallel
% %   \end{itemize}
% % \item Tagged Italian (\sh{SPEC})
% %   \\ ready to start Bulgarian, Dutch, German, Polish 
% % \end{itemize}
% % \end{frame}




% \section{Extending Wordnets}

% % \begin{frame}{Extending Wordnets}

% % \begin{itemize}
% % \item Broader
% %   \begin{itemize}
% %   \item Many new languages
% %   \item Internet slang
% %   \end{itemize}
% % \item Deeper
% %   \begin{itemize}
% %   \item Original wordnet only has Nouns, Verbs, Adjectives, Adverbs
% %     \begin{itemize}
% %     \item Added pronouns (\eng{me}), classifiers (個), 
% %       quantifiers (\eng{some}), greetings/exclamations (\eng{Hello, Thank you})
% %     \item Adding modals (\eng{must}), bound morphemes (\eng{re-, un-})
% %     \end{itemize}
% %   \item Tagging Corpora
% %     \\ how is this word used in practice?
% %   \item Richer representation of idioms
% %   \end{itemize}
% % \item Easier to build and use
% %   \begin{itemize}
% %   \item Share tools, code  and documentation
% %   \item Import information from other wordnets
% %   \end{itemize}
% % \end{itemize}
% % \end{frame}


% \begin{frame}{Many interesting things are not NVAR}
  
%   From \textit{The Adventure of the Speckled Band} (Conan Doyle,
%   1892), a short story currently being annotated as part of the NTUMC:
%   some \ul{pronouns} and \uwl{exclamatives}.
%   \begin{itemize}
%   \item \lng{``\uwl{Ah}! \ul{That} is suggestive. \uwl{Now}, on the \ul{other} side of
%     \ul{this} narrow wing runs the corridor from \udl{which} \ul{these} three
%     rooms open. \udl{There} are windows in \ul{it}, of course?''}

%   \item  \lng{``\uwl{Yes}, but very small \ul{ones}. Too narrow for \ul{anyone} to pass through.''}
    
%   \item \lng{``\uwl{Thank you}. \ul{That} is quite settled'' said \ul{he}, rising
%     and putting \ul{his} lens in \ul{his} pocket.}

    
%   \item  \lng{``\uwl{Hullo}! \ul{Here} is \ul{something} interesting''}

% \end{itemize}



% \end{frame}

% \subsection{Pronouns}


%   \begin{frame}{Pronoun Motivation and Overview}

%     \begin{itemize}
%     \item Attempting to model lexical and structural semantics
%     \begin{itemize}
%     \item For multiple languages --- identify cross-lingual differences
%     \item Exploit them to learn meaning (make the \emp{implicit} \emp{explicit})
%     \end{itemize}
%   \item Started by annotating \emp{content words} (with \emp{wordnets})
%     \item But  nouns were often translated as pronouns$_i$ 
%       --- so tag them$_i$
%       \begin{enumerate}
%       \item \emp{Identify pronouns} used in the corpus
%       \item \emp{Analyze in terms of components} --- aids matching 
%         \begin{itemize}
%         \item Extended wordnet gives \emp{full decompositional analysis}
%         \end{itemize}
%       \item \emp{Annotate} the pronouns \emp{monolingually} in each language
%         \begin{itemize}
%         \item \emp{Link to} extended wordnet for \emp{analysis}
%         \end{itemize}
%       \item \emp{Annotate} their correspondences across \emp{languages}
%       \item \emp{Analyze the distribution cross-lingually}
%       \end{enumerate}
%     \end{itemize}
%   \end{frame}
  


%     \begin{frame}{Identifying Pronouns}
%     \begin{itemize}
%        \item Examined words tagged as pronouns in (Mandarin) Chinese, English,
%          Japanese (and later Indonesian) parts of the NTU Multilingual
%          Corpus (NTU-MC) --- used the POS tags
%          \begin{itemize}
%          \item Different tag-sets identified quite different collections
%          \end{itemize}
%        \item We took the union, and filled in missing entries by hand
%          \begin{itemize}
%          \item also referred to reference grammars
%          \item not complete, but getting there 
%          \end{itemize}
%        \item 117 different types; 249 tokens:
%          %grep -v '#' extended.tab | cut -f 1 | sort -u | wc -l
%          %grep -v '#' extended.tab | cut -f 2 | grep lemma | sort | uniq -c
%          \begin{tabular}[t]{lrl}
%            Chinese & 57\\
%            English & 68\\
%            Indonesian & 40 & (in progress)\\
%            Japanese & 84 
%          \end{tabular}
%        \item We include related determiners (demonstratives and quantifiers)
%        \end{itemize}
% [numbers out of date]
%      \end{frame}
     
% % \begin{frame}{Components}%\vspace{2em}

% % \begin{itemize} \small
% % \item Type (Hypernym): \\
% %   Quantifier (no hypernym), Entity, Human, Thing, Time, Place, Manner, Reason, Quality
% % \item Person: \\
% %   1st, 2nd, 3rd, 1i, 1e
% % \item Number: \\
% %   S (singular), P (plural), D (dual)
% % \item Gender:
% %   M (Masculine), F (Feminine), N (Neuter), 
% %   GM (Grammatical Masculine), GF (Grammatical Feminine), GN (Grammatical Neuter)
% % \item Quantifier/Type: 
% %   Assertive, Elective, Negative, Other, Ever, Universal, Interrogative, Reciprocal, Reflexive, Relative
% % \item Formality (addressee honorification): Formal, Informal
% % \item Politeness (target honorification): Polite, Impolite
% % \item Proximity:
% %   Near, Distal, Medial, Remote, Close, Far
% % \item Direction: Up, Down, In, Out, From, To
% % \item Visibility: Visible, Invisible
% % \item Definiteness: Definite/Indefinite
% % \item Usage: Archaic, Colloquial
% % \end{itemize}
% % \end{frame} 

% % \begin{frame}{Analyzing Pronouns Mono-lingually}
% %        \begin{itemize}
% %        \item Decompose into: 
% %          \begin{itemize}
% %          \item \textbf{head} (\rel{hyponym})
% %          \item \textbf{restricts} (\rel{restricted-by}: new relation: 
% %            ? \rel{quantified-by})
% %          \item \textbf{features} (\rel{exemplify}: META relation)
% %          \end{itemize}
% %        \item Also mark as \rel{exemplify}$^*$ of \wn{pronoun}{n}{1} or its hyponyms
% %        \item E.g. 
% %          \begin{tabular}[t]{lll}
% %            \wn{there}{n}{1}: &\rel{hyponym} &\wn{location}{n}{1};
% %           \\ &\rel{quantified-by} &\wn{distal}{a}{1}; 
% %           \\ &\rel{exemplify} &\wn{demonstrative pronoun}{n}{1}
% %         \end{tabular}
% %         % \\  \wn{that}{a}{1}: \rel{hyponym} \wn{quantifier}{n}{1};
% %         % \\  \wn{that}{a}{1}: \rel{hyponym} \wn{quantifier}{n}{1};
% %      \end{itemize}   

% % $^*$META relation: X is an instance of Y; in theory all nouns exemplify \wn{noun}{n}{1}
% %    \end{frame}
   
% \begin{frame}{Components: Place}%\vspace{2em}
% %\begin{tikzfigure}[Head Types]
% %\MyLogo{Already refined}
% \hspace*{-3ex}\begin{tabular}{llllll}
%    Head & Type/Proximity & English & Japanese & Chinese \\
%    \hline
%  Place & Interrogative & \eng{where} &  \jpn[doko]{何処, どこ} &
%                                                                                 \cmn[nǎlǐ]{哪里}  \\
%             & Proximal      & \eng{here} & \jpn[koko]{此処, ここ} & \cmn[zhèlǐ]{这里}  \\
%             & Distal        & \eng{there}  &  & \cmn[nàli]{那里} \\
%             & ~~Medial        &   & \jpn[soko]{其処, そこ} &   \\
%             & ~~Remote        &  & \jpn[asoko]{彼処, あそこ} & \\
%             & Universal     & \eng{everywhere} & 
%                              \jpn[doko mo]{どこ も} & \cmn[dàochù]{到处}\\
%             & Existential & & \jpn[doko ka]{どこ か}
%  &  \cmn[mǒuchù]{某处} \\
%             & ~~~Assertive     & \eng{somewhere} & \\
%             & ~~~Elective     & \eng{anywhere} & \\
%             & Other         & \eng{elsewhere} & \jpn[yoso]{よそ}
%   & \cmn[biéchù]{别处} \\[2ex]
% \multicolumn{4}{l}{Not all lemmas shown}
% \end{tabular}
% \end{frame}

% \begin{frame}{Components: Place (with Korean)}
%   \tiny
% \begin{tabular}{llllll}
%    Head & Type/Proximity & English & Japanese & Chinese & Korean \\
%    \hline
%  Place & Interrogative & \eng{where} &  \jpn[doko]{何処, どこ} &
%                                                                                 \cmn[nǎlǐ]{哪里}  & \kor[eodi]{어디} \\
%             & Proximal      & \eng{here} & \jpn[koko]{此処, ここ} & \cmn[zhèlǐ]{这里}  & \kor[yeogi]{여기} \\
%             & Distal        & \eng{there}  &  & \cmn[nàli]{那里} &  \\
%             & ~~Medial        &   & \jpn[soko]{其処, そこ} &   & \kor[geogi]{거기} \\
%             & ~~Remote        &  & \jpn[asoko]{彼処, あそこ} & & \kor[jeogi]{저기} \\
%             & Universal     & \eng{everywhere} & 
%                              \jpn[doko mo]{どこ も} & \cmn[dàochù]{到处} & \kor[eodideunji]{어디든지} \\
%             & Existential & & \jpn[doko ka]{どこ か}
%  &  \cmn[mǒuchù]{某处} & \kor[eodiinga]{어딘가} \\
%             & ~~~Assertive     & \eng{somewhere} & & & \kor[eodijeongdo]{어디정도} \\
%             & ~~~Elective     & \eng{anywhere} & & & \kor[eodideunji]{어디든지} \\
%             & Other         & \eng{elsewhere} & \jpn[yoso]{よそ}
%   & \cmn[biéchù]{别处} & \kor[dareun gos]{다른 곳} \\[2ex]
% \multicolumn{5}{l}{Not checked :-)}
% \end{tabular}
% \end{frame}




% % \begin{frame}{Hypernym Examples}
% % %\MyLogo{Crude POS: a = adnominal; n = nominal; r = adverbial}
% % \begin{tabular}{lll}
% %   Type & Example & Comments \\ \hline
% %   Quantifier (no hypernym) & \lng{every, some, this$_a$} & \\ 
% %   Entity & \lng{sochira (ja)} & \\ 
% %   Human & \lng{everyone, who} & \\ 
% %   Thing & \lng{everything, what, this} & \\ 
% %   Time & \lng{anytime, when, now}& \\ 
% %   Place & \lng{anywhere, where, here}& \\ 
% %   Manner & \lng{anyhow, how}& ``in what manner''\\ 
% %   Reason &  \lng{why, hence} & ``for what reason''\\ 
% %   Quality  & \lng{konna$_a$ (ja)} & ``this kind of $\sim$''
% % \end{tabular}

% % \begin{tikzpicture}
% %     \tikzset{sibling distance=0pt} 
% %     \Tree  [.Entity Person  Thing ] 
% %   \end{tikzpicture}
% % \end{frame}


% \begin{frame}{Tagging Pronouns Mono-lingually}
%   \begin{itemize}
%   \item Tagged one document by hand \textit{The Adventure of the
%       Speckled Band}
%   \item   
%     \begin{tabular}[t]{lrrr}
%       Language & English &  Chinese & Japanese \\ \hline
%       Contentful & 1,370 & 1,177 & 463 \\
%       Other & 75 & 19 & 51\\
%       \hline 
%       Total &  1,445 & 1,196 & 514
%       \\
%       \hline 
%       Sentences & 599 & 620 & 702 \\
%       Words & 11,628 & 12,433 & 13,902 
%     \end{tabular}
%   \item Distinguished existential \lex{there} (but not dummy \lex{it})
%     with POS tags
%   \item \textit{other} includes relative pronouns, dummy
%     \lex{it}, idioms and segmentation errors
%   \end{itemize}
% \end{frame}

% \begin{frame}{Tagging and Analyzing Pronouns Cross-lingually}
%   \begin{itemize}
%   \item Automatically linked by matching features
%   \item Hand corrected:\\
%     \hspace*{-2em}
%     \begin{tabular}[t]{lrrrrrccc}
%       & \multicolumn{6}{c}{Linked Pronouns} & \multicolumn{2}{c}{Non-linked Pronouns} \\
%       &\multicolumn{5}{c}{\# Matching Features} & Pronoun & English & Other \\
%       & 5 & 6 & 7 & 8 & 9 & to Noun & & \\
%       \hline
%       \# Chinese   & 5 & 19 & 54 & 789 & 58 & 134 & 369 & 215 \\
%       \# Japanese  & 15 & 120 & 114 & 37 & 32 & 139 & 943 & 109
%     \end{tabular}
%   \item Case and politeness mismatches common
%   \item A surprising number of non-linked pronouns in Chinese and Japanese
%   \end{itemize}
% \end{frame}


% \begin{frame}{Interesting Cross-Linguistics Differences}
%   \begin{exe}
%     \ex \wemp{She}$_i$ shot \wemp{him}$_j$ and then \wemp{herself}$_i$
%     \begin{xlist}
%       \ex \gll \wemp{奥-さん} が \wemp{旦那-さん} を 撃って 、 それから \wemp{自分} も 撃った \\
%       oku-san ga danna-san wo utte , sorekara jibun mo utta \\
%       % wife-HON NOM husband-HON ACC shoot-CONJ , and+then self too shoo-PST \\
%       \trans \wemp{Wife}$_i$ shot \wemp{husband}$_j$ and then shot \wemp{self}$_i$ too
      
%       \ex  \gll \wemp{她} 拿 枪  先 打 \wemp{丈夫} , 然后 打 \wemp{自己}\\
%       t\={a} ná qi\={a}ng xi\={a}n d\v{a} zhàngf\={u} , ránh\`{o}u d\v{a} zìj\v{\i} \\
%       % 3SG take gun first shoot husband , and+then shoot self\\
%       \trans \wemp{She}$_i$ took the gun to first shoot \wemp{husband}$_j$, and then shot \wemp{self}$_i$
%     \end{xlist}
%   \end{exe}
% \end{frame}

% \begin{frame}{Interesting Cross-Linguistic Differences}
% \begin{exe}
% \ex\label{s:none}    \mbox{[many (cases) strange] \ldots but \wemp{none} commonplace \ldots}
% \begin{xlist}
%   %\CJKfamily{gbsn}
%   \ex \gll      但是    却   \wemp{没有}    \wemp{一例} 是  平淡无奇      的 \\
%   Dan4shi4 que4 mei2you3 yi1li4 shi4 ping2dan4wu2qi2 de \\
%   %\CJKfamily{min}
%   \trans ‘But, there is \wemp{not  one case} that is featureless.’
%   \ex \gll  \wemp{どれ も} 尋常で は \wemp{ない} 事件 である \\
%   \wemp{Dore mo}   jinjode   wa \wemp{nai}    jiken dearu \\
%   \trans ‘\wemp{Everything} is a case which is \wemp{not} usual.’
% \end{xlist}
%   \ex\label{s:dif}  \wemp{It} is a swamp adder!
%   \begin{xlist}
%     %\CJKfamily{gbsn}
%     \ex \gll \wemp{这}     是   一 条      沼地        蝰蛇 ！ \\
%     Zhe4 shi4 yi1tiao2 zhao3di4 kui2she2 ! \\
%     %\CJKfamily{min}
%     \trans `\wemp{This} is a swamp adder!'
%     \ex \gll  沼蛇 だ ！ \\
%     numahebi da ! \\
%     \trans `\wemp{$\phi$} is a swamp snake'
%   \end{xlist}
% \end{exe}
% \end{frame}



% \subsection{Interjections}
% \begin{frame}{Interjections}
%   \begin{itemize}
%   \item words or phrases that
%     \begin{itemize}
%     \item constitute a whole linguistic act
%       \\ (do not combine in integrated syntactic constructions)
%     \item do not refer to events, but instead carry expressive meaning
%       \\ (Huddleston and Pullum, 2002)
%     \end{itemize}
%   \item     We follow Jovanovíc (2004) and Ameka (1999), and use the term broadly,
%     covering plain interjections, greetings and many more \ldots
%   \item We introduce a new POS \textbf{x} for these non-referential expressions
%     (also used for numeral classifiers in languages like Japanese)
%     \begin{itemize}
%     \item definition with the form ``an expression that is uttered ...:
%     \item \rel{exemplifies}: utterance (07109847-n)
%     \item enrich this flatter hierarchy with links to other existing
%       concepts (when possible)
%     \end{itemize}
%   \end{itemize}
% \end{frame}

% \begin{frame}{What does our broad sense of interjections include?}
%   \begin{itemize}
%   \item expressions of emotion, such as surprise, disgust, etc.
% \\(e.g. \lng{wow, ugh, yuk, gosh, ...})
% \item expressions used in greetings, leave-taking, thanking, apologizing, etc.
% \\(e.g. \lng{hello, thank you, goodbye, ...})
% \item expressions used for swearing
% \\ (e.g. \lng{damn, shit, bite me, ...})
% \item expressions used in responding
% \\(e.g. \lng{yes, no, OK, yeah, you bet, ...})
% \item and a long tail of onomatopoeia
%   \\(e.g. \lng{ding-dong, woof-woof, miao, ...})
% \end{itemize}
% \end{frame}

% \begin{frame}{New Interjective Senses}

%   \begin{center}\small
%     \hspace*{-0.75em}\begin{tabular}{ l r l r}
%       Concept & Senses  & Concept & Senses \\ \hline
%       Surprise, Wonder & 58 & Pity, Sorrow & 19 \\ 
%       Joy, Pleasure & 17 & Anger, Annoyance, Irritation & 41 \\ 
%       Approval, Triumph, Enthusiasm & 10 & Contempt, Disgust, Impatience & 59 \\
%       Pain & 7 & Sympathy & 2 \\ 
%       Delight & 11 & Fear & 3 \\ 
%       Relief & 2 & Encouragement & 16 \\ 
%       Attention-Seeking & 36 & Toasting & 10 \\ 
%       General Greetings & 13 & Morning Greetings & 2 \\ 
%       Afternoon Greetings & 2 & Night Greetings & 2 \\ 
%       General Farewells & 21 & Night Farewells & 5 \\ 
%       Checkmate & 2 & & \\ 
%       \hline
%       \textbf{Total number of senses} & 336 & & \\ 
%     \end{tabular}
%   \end{center}
% \end{frame}

% \begin{frame}{New Interjective Senses}

%   \begin{center}
% \begin{avm}
% \[ \wnida{\wnid{80000001-x} (general greeting)}
% lemmas & \parbox[t]{24em}{\lex{aloha, ciao, g'day, good day, hallo, halloa, halloo, hallow, hello, hi, howdy, hullo, 'sup}} \\ 
% definition & \parbox[t]{24em}{an expression that is uttered as a general greeting, regardless of the time of day} \\ 
% exemplifies & \wnid{15167474-n} (utterance) \\ 
% see also & \wnid{06630017-n} (greeting) \]
% \end{avm} \\

% \begin{avm}
% \[ \wnida{\wnid{80000002-x} (checkmate)}
% lemmas &\lex{checkmate, mate} \\ 
% definition & \parbox[t]{24em}{an expression that is uttered during a game of chess to declare that the final winning move has taken place} \\ 
% exemplifies & \wnid{15167474-n} (utterance) \\ 
% see also & \wnid{00167764-n} (checkmate) \]
% \end{avm} \\
% \end{center}
% \end{frame}


% \subsection{Metaphors and Metonymy}
% \label{sec:metaphors-metonymy}

% \begin{frame}{Metaphors and Metonymy}
%   \begin{itemize}
%   \item We want to know how different senses of a word relate to each other 
%   \item So we try to link different senses of the same word
%     \begin{itemize}
%     \item \txx{metonymy} referring to a thing by its attribute
%     \item \txx{metaphor}  a non-literal comparison between two unlike things .
%     \end{itemize}
%   \begin{center}
% \begin{tikzpicture}[derivation]
% \node[lit_definition, text width=2.7cm] (1) at (0,0) {\sense{march}{{2}} [\synonym{marching}] the act of marching; walking with regular steps};
% \node[ass_definition, text width=2.cm] (3) [right = 1.6cm of 1] {\sense{march}{{4}} a procession of people walking together};
% \node[lit_definition, text width=2.9cm] (4) [below right = .6cm and -2.8cm of 1] {\sense{march}{{1}} the month following February and preceding April};

%     \path[related] (1) edge node[anchor=south, font={\footnotesize}] {\relatededgelabel} (3);
% \node[met_definition, text width=2.6cm] (2) [below = .65cm of 3] {\sense{march}{{3}} a steady advance, e.g.\ ``the march of science''};
%         \path[analogy] (3) edge node[anchor=west, font={\footnotesize}] {\metaphoricaledgelabel} (2);
%       \end{tikzpicture}
% \end{center}      
% \item Unlinked senses will by \txx{homonyms}
% \end{itemize}

% \end{frame}

% \begin{frame}{Work in Progress}
%   \begin{itemize}
%   \item Work done mainly by Rowan Hall Maudslay  \citep{maudslay-etal-2024-chainnet-structured}
%   \item Links senses (word + synset combinations) as metaphor or metonymy
%   \item English Wordnet (Princeton 3.0) annotated for almost all words
%     with 3 or more senses
%   \item 6,500 entries annotated (almost all words with three or more senses)
%   \item Plan to annotate the rest after Rowan finishes his PhD
%   \end{itemize}

% \end{frame}


% \begin{frame}{We can see how English metaphors are used cross-lingually}
%   \begin{itemize}
%   \item For example, translations of linked pairs likely to be
%     translated with the same word
%     \begin{itemize}
%     \item So  metaphor/metonomy holds cross language
%     \item (OR the wordnet was built automatically and badly)
%     \end{itemize}
%   \item For example \eng{head} ``body part'' is metaphorically
%     extended to ``person in charge'' in English and this is also true in Japanese,
%     頭, Italian, \eng{capo} and many other languages
%   \item But not all metaphors are shared  \eng{head} ``main part of a grammatical constiuent'' is not 頭 in Japanese
    

%   \end{itemize}
% \end{frame}

% \begin{frame}{Measure the translation overlap with OMW}

%   \begin{itemize}
%   \item For every English word that has been annotated
%     \begin{itemize}
%     \item For each pair of English senses
%     \item Look up all translations of each sense
%       \\ measure the overlap (we use jaccard)
%       \item Store and note if the sense are linked or not
%     \end{itemize}
%   \item So for example \eng{head} has 3 senses (really many more)
%     \begin{enumerate}
%     \item head, caput
%       \\ ヘッド, \ul{頭}, 頭部 
%     \item head, chief, top dog
%       \\ 大頭, 主任者, 御頭, 頭領, \ul{頭}
%     \item head,  head word
%       \\ 主要語
%     \end{enumerate}
%   \item Overlap
%     \begin{itemize}
%     \item 1-2: 0.125 %(/ 1.0  8) 0.125
%     \item 2-3: 0
%     \item 1-3: 0
%     \end{itemize}
%   \end{itemize}
  

% \end{frame}

% \begin{frame}[allowframebreaks]{We do this for all the lexicons in OMW-1.4}

 
%   \begin{tabular}{lrrrr}
%     Lang & Unlinked & Metaphor & Metonomy & Non-zero \\ \hline
% sv & 0.70 & 1.71 & 2.10 & 417   \\
% he & 0.69 & 1.58 & 2.31 & 472   \\
% eu & 0.82 & 1.30 & 1.80 & 5562  \\
% ja & 0.89 & 1.11 & 1.55 & 8130 \\
% ro & 0.83 & 1.30 & 1.71 & 6429 \\
% sq & 0.74 & 1.27 & 2.38 & 822 \\
% en & 0.97 & 1.08 & 1.07 & 33415 \\
% da & 0.74 & 1.46 & 2.09 & 332 \\
% lt & 0.63 & 1.46 & 2.86 & 692 \\
% es & 0.85 & 1.17 & 1.79 & 5915 \\
% hr & 0.81 & 1.22 & 1.98 & 3576 \\
%     fr & 0.94 & 1.10 & 1.25 & 18975 \\
%   \end{tabular}

%   Divide the distance for each group by the average distance for all senses.

%   \begin{tabular}{lrrrr}
%     Lang & Unlinked & Metaphor & Metonomy & Non-zero \\ \hline
% nb & 0.73 & 1.49 & 2.12 & 339 \\
% bg & 0.73 & 1.54 & 2.11 & 380 \\
% gl & 0.70 & 2.13 & 1.55 & 271 \\
% cmn & 0.74 & 1.57 & 1.97 & 1347 \\
% it & 0.80 & 1.39 & 1.81 & 4614 \\
% arb & 0.74 & 1.28 & 2.38 & 1337 \\
% pl & 0.75 & 1.43 & 2.07 & 1326 \\
% el & 0.72 & 1.30 & 2.45 & 1274 \\
% iwn & 0.79 & 1.46 & 1.80 & 968 \\
% nn & 0.73 & 1.52 & 2.10 & 340 \\
% id & 0.87 & 1.21 & 1.58 & 9803 \\
%   \end{tabular}
%   \begin{tabular}{lrrrr}
%     Lang & Unlinked & Metaphor & Metonomy & Non-zero \\ \hline
% is & 0.73 & 1.51 & 2.13 & 388 \\
% ca & 0.83 & 1.23 & 1.84 & 6000 \\
% sk & 0.73 & 1.35 & 2.31 & 2067 \\
% fi & 0.82 & 1.35 & 1.76 & 7971 \\
% nl & 0.79 & 1.38 & 1.90 & 2991 \\
% sl & 0.81 & 1.41 & 1.74 & 6243 \\
% th & 0.80 & 1.48 & 1.71 & 2336 \\
% pt & 0.80 & 1.26 & 1.96 & 5414 \\
% zsm & 0.87 & 1.22 & 1.60 & 10272 \\
%   \end{tabular}
  
%   Unlinked  share fewer translations, metaphor shares more and metonomy shares much more.

%   \Large Metonomy is more likely to hold cross linguistically!
  
% \end{frame}

% \begin{frame}{Can we use other lemmas to classify links}

%   \begin{itemize}
%   \item Synsets have sets of synonyms, so the other senses may include useful information
%   \item Consider the synsets with a sense of chestnut
%     \begin{itemize}
%     \item[s1] chestnut ``edible nut of any of various chestnut trees of the genus Castanea'' (NATURAL\_BODY)

%     \item[s2] chestnut, chestnut tree ``any of several attractive deciduous trees yellow-brown in autumn; yield a hard wood and edible nuts in a prickly bur'' (PLANT)
%     \item[s3] chestnut ``wood of any of various chestnut trees of the genus Castanea''
%       (CHEMICAL!)
%     \end{itemize}
%   \item These are annotated as s1 → s2 → s3 (metonomy)
%   \item If we look at all the senses
%     \begin{itemize}
%     \item  s1 → s2 has a diff: ``+ tree''
%     \item s2 → s3 has a diff: ``- tree''
%     \end{itemize}
%   \item These patterns repeat \ldots
%     \\ Cherry → Cherry Tree → Cherry\hfill (en)
%   \end{itemize}

 
% \end{frame}
% \begin{frame}{They can be found across languages}

%   \begin{itemize}
%   \item s1 → s2: marron →	marronnier \hfill (fr)
%   \item s1 → s2: kastanja →	kastanjapuu \hfill (fi)
%   \item  s1 → s2:	kastanje →	kastanjeboom \hfill (nl)
%   \item  s1 → s2:	castană →	castan \hfill (ro)
%   \item  s1 → s2:	castanăcomestibilă →	castan \hfill (ro)
%   \end{itemize}	
% \end{frame}


% \begin{frame}{Summary}
  
% \begin{itemize}
% \item The combination of detailed annotation of meaning and
%   cross-lingual links allow us to explore differences in meaning
%   representation
% \item We will add to wordnet and make searchable
% \item Richer wordnets in more languages allow us to investigate more phenomena
% \item We have not yet looked at use in sense-tagged corpora
% \item LLMs may capture these regularities, but are hard to explore with
%   \begin{itemize}
%   \item They tend to lose the longs tail
%   \item Colleagues in Denmark are seeing if they capture Danish specific metaphors
%   \end{itemize}
% \end{itemize}

% \end{frame}


% \section{CILI: the Collaborative InterLingual Index}
% \begin{frame}{Why the Collaborative InterLingual Index?}
% We want to make it easier to link things together.
%   \begin{itemize}
%   \item There are wordnets for many languages
%     \begin{itemize}
%     \item Currently they link through PWN (3.0)
%     \end{itemize}
%   \item Many projects are adding new synsets 
%     \begin{itemize}
%     \item And not just synsets: lemmas, relations, POS, meta-data
%       (domains, sentiment \ldots)
%     \end{itemize}
%   \item We want to be able to link them even if they are not in PWN
%   \item We want to minimize wasted effort
%     \begin{itemize}
%     \item Adding the same thing in different projects
%     \item Fixing the same errors in different projects
%     \end{itemize}
%   \item We want to spread the burden of development
%   \end{itemize}
  
% \end{frame}



% \begin{frame}{The Princeton Wordnet of English}

%   \begin{center}
% \begin{tikzpicture}
% \tikzset{level distance=12ex}
% \Tree [.\node(ea){\begin{avm}% 
%     \[ \wnida{pwn:0121-n}
%     lemma & \eng{sibling} \\ 
%     \]\end{avm}};
%       [.\node(eb){\begin{avm}% 
%     \[ \wnida{pwn:0124-n}
%     lemma & \eng{brother} \\ 
%     \]\end{avm}}; ]
%       [.\node(ec){\begin{avm}% 
%     \[ \wnida{pwn:0127-n}
%     lemma & \eng{sister} \\ 
%      \]\end{avm}}; ]
%          ]
%     \end{tikzpicture}
%   \end{center}
%     \begin{itemize}
%     \item Models the lexical structure of English
%     \end{itemize}

%   \end{frame} 
  
% \begin{frame}{Japanese Wordnet (transfer structure) ($-$)}

%   \begin{center}
% \begin{tikzpicture}
% \tikzset{level distance=12ex}
% \Tree [.\node(ja){\begin{avm}% 
%     \[ \wnida{pwn:0121-n}
%     lemma & \eng{姉妹兄弟, 兄弟} \\ 
%     \]\end{avm}};
%       [.\node(jb){\begin{avm}% 
%     \[ \wnida{pwn:0124-n}
%     lemma & \eng{兄弟, 男兄弟} \\ 
%     \]\end{avm}}; ]
%       [.\node(jc){\begin{avm}% 
%     \[ \wnida{pwn:0127-n}
%     lemma & \eng{姉妹} \\ 
%     \]\end{avm}}; ]
%          ]
%     \end{tikzpicture}
%   \end{center}
%     \begin{itemize}
%     \item Saves time by copying the English structure
%     \item But misses some concepts (younger and older siblings)
%     \end{itemize}


%   \end{frame}
  
% \begin{frame}{Japanese Wordnet (transfer structure) ($+$)}

%   \begin{center}
% \begin{tikzpicture}
% \tikzset{level distance=12ex}
% \Tree [.\node(ja){\begin{avm}% 
%     \[ \wnida{pwn:0121-n}
%     lemma & \eng{姉妹兄弟, 兄弟} \\ 
%     \]\end{avm}};
%       [.\node(jb){\begin{avm}% 
%     \[ \wnida{pwn:0124-n}
%     lemma & \eng{兄弟, 男兄弟, \ul{兄}, \ul{弟}} \\ 
%     \]\end{avm}}; ]
%       [.\node(jc){\begin{avm}% 
%     \[ \wnida{pwn:0127-n}
%     lemma & \eng{姉妹, \ul{姉}, \ul{妹}} \\ 
%     \]\end{avm}}; ]
%          ]
%     \end{tikzpicture}

%   \end{center}\begin{itemize}
%     \item Saves time by copying the English structure
%     \item But merges dis-similar concepts (younger and older brother)
%     \end{itemize}

% \end{frame}


% \begin{frame}{OMW 1.0}


%   \begin{center}
% \begin{tikzpicture}
% \tikzset{level distance=15ex}
% \Tree [.\node(ja){\begin{avm}% 
%     \[ \wnida{pwn:0121-n}
%     lemma-eng & \eng{sibling} \\ 
%     lemma-jpn & \eng{姉妹兄弟} \\ 
%     \]\end{avm}};
%       [.\node(jb){\begin{avm}% 
%     \[ \wnida{pwn:0124-n}
%     lemma-eng & \eng{brother} \\ 
%     lemma-jpn & \eng{兄弟} \\ 
%     \]\end{avm}}; ]
%       [.\node(jc){\begin{avm}% 
%     \[ \wnida{pwn:0127-n}
%     lemma-eng & \eng{sister} \\ 
%     lemma-jpn & \eng{姉妹} \\ 
%     \]\end{avm}}; ]
%          ]
%     \end{tikzpicture}
%   \end{center}
%     \begin{itemize}
%     \item Uses the English structure 
%     \item Loses any nodes not linked to English
%     \end{itemize}
  
%   \end{frame}
  
% \begin{frame}{Japanese Wordnet (native structure)}

% \begin{tikzpicture}
% \tikzset{level distance=12ex}
% \Tree [.\node(ja){\begin{avm}% 
%     \[ \wnida{jwn:3478-n}
%     lemma & \eng{姉妹兄弟, 兄弟} \\ 
%     \]\end{avm}};
%       [.\node(jb){\begin{avm}% 
%     \[ \wnida{jwn:6597-n}
%     lemma & \eng{兄弟, 男兄弟} \\ 
%     \]\end{avm}}; 
%  [.\node(je){\begin{avm}% 
%     \[ \wnida{jwn:6897-n}
%     lemma & \eng{兄} \\ 
%     \]\end{avm}}; ]
%  [.\node(jf){\begin{avm}% 
%     \[ \wnida{jwn:5957-n}
%     lemma & \eng{弟} \\ 
%     \]\end{avm}}; ]
% ]
%       [.\node(jc){\begin{avm}% 
%     \[ \wnida{jwn:8745-n}
%     lemma & \eng{姉妹} \\ 
%     \]\end{avm}};  [.\node(jg){\begin{avm}% 
%     \[ \wnida{jwn:2897-n}
%     lemma & \eng{姉} \\ 
%     \]\end{avm}}; ]
%  [.\node(jh){\begin{avm}% 
%     \[ \wnida{jwn:2957-n}
%     lemma & \eng{妹} \\ 
%     \]\end{avm}}; ]]
%          ]
%     \end{tikzpicture}

%     \begin{itemize}
%     \item More accurately models Japanese
%     \end{itemize}

%   \end{frame}
  
% \begin{frame}{OMW 2.0}


% \hspace*{-1em}\begin{tikzpicture}
% \tikzset{level distance=15ex}
% \Tree [.\node(ja){\begin{avm}% 
%     \[ \wnida{cili:i1234 = pwn:0121-n, jwn:3478-n}
%     lemma-eng & \eng{sibling} \\ 
%     lemma-jpn & \eng{姉妹兄弟, 兄弟} \\ 
%     \]\end{avm}};
%       [.\node(jb){\begin{avm}% 
%     \[ \wnida{cili:i1245 = pwn:0124-n, jwn:6597-n}
%     lemma-eng & \eng{brother} \\ 
%     lemma-jpn & \eng{兄弟, 男兄弟} \\ 
%     \]\end{avm}}; [.\node(je){\begin{avm}% 
%     \[ \wnida{jwn:6897-n = cili:?}
%     lemma-jpn & \eng{兄} \\ 
%     \]\end{avm}}; ]
%  [.\node(jf){\begin{avm}% 
%     \[ \wnida{jwn:5957-n = cili:?}
%     lemma-jpn & \eng{弟} \\ 
%     \]\end{avm}}; ] ]
%       [.\node(jc){\begin{avm}% 
%     \[ \wnida{cili:i7643 = pwn:0127-n, jwn:8745-n}
%     lemma-eng & \eng{sister} \\ 
%     lemma-jpn & \eng{姉妹} \\ 
%     \]\end{avm}};   [.\node(jg){\begin{avm}% 
%     \[ \wnida{jwn:2897-n = cili:?}
%     lemma-jpn & \eng{姉} \\ 
%     \]\end{avm}}; ]
%  [.\node(jh){\begin{avm}% 
%     \[ \wnida{jwn:2957-n = cili:?}
%     lemma-jpn & \eng{妹} \\ 
%     \]\end{avm}}; ]]
%          ]
%     \end{tikzpicture}

%     \begin{itemize}
%     \item Unifies the structures of all languages
%     \item Linked by the Collaborative Interlingual Index (CILI)
%     \item Requires some coordination, \ldots
%     \end{itemize}
%   \end{frame}



% \begin{frame}{The solution: an InterLingual Index (ILI)}
% \begin{enumerate}
% \item The Interlinear Index (ILI) should be a flat list of concepts (and instances).
% \item\label{ili:links} The semantic and lexical relations should mean the same things for all languages.
% \item Concepts should be constructed for salient and frequent lexicalized
%   concepts in all languages.
% \item Concepts linked to Multiword units (MWUs) in wordnets should be included.  
% %As space is not an issue, the default should be to include concepts based on common MWUs in languages.
% \item A formal ontology could be linked to but separate from the wordnets.
% \end{enumerate}

% Basic idea well known (Fellbaum and Vossen, 2008). 
% \end{frame}

% \begin{frame}{Collaboratively Developed (CILI) \ldots}
% \begin{enumerate}
% \item\label{license} The license must allow redistribution of the index
% \item\label{persistent} ILI IDs should be persistent:
%  we never delete, only \txx{deprecate} or \txx{supercede};
%   we should not change the meaning of the concept
% \item\label{defs} Each new ILI concept should have a definition in English, as
%  this is the only way we can coordinate across languages.
% \\
%   The definition should be unique (which is not currently true) and 
%   preferably also parse and sense tag information should be included.
%   Definition changes will be moderated.
% \end{enumerate}
% \end{frame}

% \begin{frame}{\ldots Collaboratively Developed (CILI)}
% \begin{enumerate}
% \item\label{prov} Each new ILI concept should link to a synset in an existing project that is part of the OMW with one of a set of known relations (\rel{hypernymy}, \rel{meronomy}, \rel{antonymy}, \ldots) 
% \item\label{link} This synset should link to another synset in an existing project that is part of the OMW and links to an ILI concept.
%   \begin{itemize}
%   \item[$\Rightarrow$]  each concept is linked to another concept through at least one wordnet in the grid
%   \end{itemize}
%   % \item[?]\label{exes} Each new synset should have an example
% \item\label{add} Any project adding new synsets should first check
%   that they do not already exist in the CILI
%   \begin{itemize}
%   \item New concepts are added through their existing in a wordnet
%   \item If something fulfills the criteria is proposed
%   \item If no objections after three months then it is added
%   \end{itemize}
% \end{enumerate}

% \end{frame}

% \section{The Open Multilingual Wordnet 2.0}

% \begin{frame}{Open Multilingual Wordnet 2.0}
% \begin{itemize}
% \item OMW and merging hosted at UP, proxied through omwn
%   \\ \url{http://omwn/omw/}
%   \\ show provenance and confidence
% \item Individual wordnets can be uploaded
% \item ILI as RDF --- shared on github
% \item Each new wordnet release can suggest ILI candidates
%   \begin{itemize}
%   \item Mark as \texttt{ili='in'} ``ILI new''
%   \item Only projects can add new ILI entries
%   \end{itemize}
% \item Review period (1--3 months)? --- accept if no comments
% \item Allow orthographic variants
% \item Documentation (and all code) hosted on github
% \end{itemize}
% \end{frame}

% % \begin{verbatim}
% % @prefix ili: <http://ili.org/> .
% % @prefix ns1: <http://purl.org/dc/elements/1.1/> .
% % @prefix ns2: <http://www.w3.org/2004/02/skos/core#> .
% % @prefix ns3: <http://www.w3.org/2002/07/owl#> .
% % @prefix pwn30: <http://pwn30.org/> .
% % @prefix rdf: <http://www.w3.org/1999/02/22-rdf-syntax-ns#> .
% % @prefix rdfs: <http://www.w3.org/2000/01/rdf-schema#> .
% % @prefix xml: <http://www.w3.org/XML/1998/namespace> .
% % @prefix xsd: <http://www.w3.org/2001/XMLSchema#> .

% % ili:ili1 a ns3:concept ;
% %     ns1:source pwn30:s00001740-a ;
% %     ns3:sameAs pwn30:s00001740-a ;
% %     ns2:definition "(usually followed by `to') having the necessary means or skill or know-how or authority to do something"@en .

% % ili:ili10 a ns3:concept ;
% %     ns1:source pwn30:s00003553-s ;
% %     ns3:sameAs pwn30:s00003553-s ;
% %     ns2:definition "coming into existence"@en .
% % \end{verbatim}
% %\end{frame}


% %%
% %% PWN grandfathering
% %%

% % \begin{frame}{Grandfathering}

% % \begin{itemize}
% % \item Start off with Princeton Wordnet 3.0 Synsets
% % \\ even if they have
% % \begin{itemize}
% %   \item Non-unique definitions   
% %     \\ almost fixed
% %   \item Ill-formed definitions (`;', `,'), too short 
% %     \\ short definitions need to be revised
% %     \\ very hard for e.g. quantifiers
% %   \item Not linked to another synset (orphans)
% %     \\ any volunteers?
% %   \end{itemize}
% % \end{itemize}

% % \end{frame}


% % \begin{frame}{LMF}
% %   \begin{itemize}
% %   \item add an 'ili' attribute to synsets
% %   \item encourage more META data
% %   \end{itemize}

% % \end{frame}
% \section{Future Work}
% \begin{frame}{Where do we go from here?}
%   \begin{itemize}
%   \item Convert existing wordnets to LMF
%     \begin{itemize}
%     \item \Large \txx{Validate them}
%     \end{itemize}
%     \small Possibly revise the LMF as needed
%   \item Upload it to the Grid
%     \begin{itemize}
%     \item We will validate it again
%     \item We add the wordnet to OMW: linking through ILI
%     \item Look for \texttt{ILI='in'} ``ILI new''
%     \item Check the definition is good (and not too close); parse it?
%     \item Check it is linked to something existing
%     \end{itemize}
%   \item Finish New OMW interface
%   \item Add an interface for changing definitions, deprecating and
%     superseding
%   \end{itemize}
% \end{frame}


% \begin{frame}{Where do we go from here?}
%   \begin{itemize}
%   \item New POS: 'x', 'q', 'c', 'p', \ldots
%   \item Orthographic variants explicitly modeled
%   \item More languages, links, words
%   \item Sense tagged corpora and live examples
%   \item More sense groupings  (iliXXX $\approx$ iliYYY) 
%   \item Decompositional semantics: \lex{un-clasp}, \lex{today} ``this day'' 
%   \item External DBS: \texttt{species, geo-wordnet, images, National
%       Cancer Institute Thesaurus, \ldots}
%   \item Integration with vector-space representations
%     \begin{itemize}
%     \item modeling fuzziness
%     \item sense-embeddings
%     \item typing relations
%     \end{itemize}
%   \end{itemize}
% \end{frame}

% \begin{frame}{What will this enable?}

%   \begin{itemize}
%   \item Better WSD
%     \begin{itemize}
%     \item Fewer indistinguishable senses
%     \item More training data
%     \item More relations
%     \end{itemize}
%   \item New synsets and senses:
%     \begin{itemize}
%     \item \wn{smart phone}{n}{1} ``a mobile phone with more advanced computing capability and connectivity than basic feature phones''
%     \item \wn{klunen}{v}{1} ``to walk on skates across non-ice'' (nl)
%     \item \wn{satay}{n}{1} ``A popular dish made from small pieces of meat or fish grilled on a skewer and served with a spicy peanut sauce (Nusuntara)''
%     \item around 20-30,000 coming soon
%     \item[\ldots] and now it is easier to add them!
%     \end{itemize}
%   \end{itemize}
% \end{frame}

\section{Wordnets}
\begin{frame}
\frametitle{Introduction to Wordnets}

\begin{itemize}
    \item A wordnet is a lexical database that organizes words into sets of synonyms, called synsets.
    \item Each synset represents a unique concept, capturing different senses of a word.
    \item Synsets are connected by semantic relations like hypernyms (superordinate) and hyponyms (subordinate).
    \item Wordnets are used in linguistic research, natural language processing (NLP), artificial intelligence, education and (a bit) language maintenance.
\end{itemize}
\end{frame}

\begin{frame}
\frametitle{History of Wordnets}

\begin{itemize}
\item Wordnet was first developed for English at Princeton University \citep{miller1995wordnet}.
  
    \item It was designed it as a model of human lexical memory.
    \item It became a foundational tool for computational linguistics and NLP.
    \item It has influenced the development of lexical resources for many other languages.
    \item Currently the Open English WordNet is maintained by John
      McCrae.  It is extending and revising the Princeton Wordnet.
\end{itemize}
\end{frame}

\begin{frame}{Wordnet Relations}
  \begin{itemize}
  \item Constitutive
    \\  \wn{driver}{n}{1} \rel{hyponym} \wn{operator}{n}{1}
    \\  \wn{driver}{n}{1} \rel{antonym} \wn{nondriver}{n}{1}
    \\   \wn{driver}{n}{3} \rel{synonym} \wn{number one wood}{n}{1}	  
   \\   \wn{golf club}{n}{1} \rel{meronym} \wn{clubhead}{n}{1}	  
  \item Derivational
    \\  \wn{driver}{n}{1} \rel{derived} \wn{drive}{v}{1} 
  \item Descriptive
    \\ \wn{driver}{n}{2} \rel{domain} \wn{computer science}{n}{1}
  \item Implied (from e.g. definitions)
    \\ \wn{driver}{n}{1} ``the \wn{operator}{n}{1} of a \wn{motor
      vehicle}{n}{1}''
    \\ \wn{driver}{n}{2} ``a \wn{program}{n}{3} that
    \wn{determines}{v}{2} how a \wn{computer}{n}{2} will
    \wn{communicate}{n}{1} with a \wn{peripheral device}{n}{1}''
  \end{itemize}
  Over 30 relations are defined, and more are being added \\
  (such as \rel{metaphor}, \rel{metonymy}, \ldots)
\end{frame}


\begin{frame}{The synset for  \eng[tiger]{虎}}
 \includegraphics[width=0.9\textwidth]{pics/wn-jpn-tiger.pdf}
\end{frame}

\begin{frame}
\frametitle{Wordnet as a Graph}

\begin{itemize}
    \item Synsets can be represented as nodes, and semantic relations as edges.
    \item Hypernymy links form a hierarchical tree, with more general concepts at higher levels.
      \begin{itemize}
      \item Actually a directed acyclic graph, as multiple inheritance is possible
      \end{itemize}
    \item Wordnets also include lateral connections like meronyms and antonyms.
    \item The graph structure allows for efficient computation of semantic similarity.
    \item Graph-based algorithms are used in tasks like word sense disambiguation and knowledge extraction.
\end{itemize}
\end{frame}



\begin{frame}{Usage}
  \begin{itemize}
  \item Wordnets are used for 
    \begin{itemize}
    \item studying word similarity 
    \item describing the meaning of texts
    \item testing WSD algorithms
    \item improving parse ranking
    \item machine translation (not so much these days)
    \item puzzle generation and solving
    \item joke generation
    \item distinguishing senses for word embeddings
    \item language learning
    \end{itemize}
  \item They are still useful in the age of LLMs
    \begin{itemize}
    \item As a check on quality
    \item For explicit reasoning
    \item For languages with insufficient language data
    \item For teaching
    \end{itemize}
  \end{itemize}

\end{frame}

\begin{frame}
\frametitle{Multilingual Wordnets}

\begin{itemize}
    \item Wordnets exist for over 100 languages, adapted from the original Princeton WordNet model.
    \item Multilingual wordnets link synsets across different languages, mapping equivalent concepts.
    \item Cross-lingual mappings help in tasks like machine translation and multilingual semantic analysis.
    \item The Global WordNet Association (GWA) promotes the development of wordnets in multiple languages.
    \item Projects like the Open Multilingual Wordnet (OMW) provide free access to linked wordnets.
\end{itemize}
\end{frame}

\begin{frame}
\frametitle{Challenges in Building Wordnets}

\begin{itemize}
    \item Developing wordnets requires careful analysis which cannot be fully automated.
    \item Word meaning can be culturally specific, requiring careful consideration of local concepts.
    \item Data collection for under-resourced languages often involves fieldwork and expert collaboration.
    \item Validation and refinement of synsets is an ongoing process, especially for evolving languages.
    \item Many language resources are not open, so data cannot easily be shared.
\end{itemize}
\end{frame}

% \begin{frame}
% \frametitle{Building Wordnets for Under-resourced Languages}

% \begin{itemize}
%     \item Under-resourced languages may lack dictionaries, corpora, or standardized orthographies.
%     \item Wordnets for these languages require extensive collaboration with native speakers.
%     \item Data collection can involve workshops, interviews, and domain-specific lexicons.
%     \item Wordnet development can support language preservation and revitalization efforts.
%     \item These wordnets help document endangered languages and make them accessible for future research.
% \end{itemize}
% \end{frame}

\begin{frame}
\frametitle{Open Multilingual Wordnet (OMW)}

\begin{itemize}
    \item OMW is a free multilingual lexical database linking wordnets from many languages \citep{Bond:Foster:2013}
    \item OMW includes wordnets for widely spoken languages like English, Spanish, and Chinese
    \item It also contains wordnets for under-resourced and endangered languages
    \item Researchers can access and contribute to OMW to support cross-linguistic studies
    \item The platform promotes collaborative development and data sharing across language communities
    \item In particular it has pushed open-licenses for the wordnets
\end{itemize}
\end{frame}

\begin{frame}{Wordnets in the world 2008}
\begin{center}
   \noindent\includegraphics[width=1.0\textwidth]{pics/wns-2008}
\end{center}
Green is free; Blue is research only; Brown costs money
\end{frame}

\begin{frame}{Wordnets in the world 2023}
\begin{center}
   \noindent\includegraphics[width=1.0\textwidth]{pics/wns-2023}
\end{center}
Green is free; Blue is research only; Brown costs money
\end{frame}

\begin{frame}{The situation is much better!  But not perfect}
  
\begin{itemize}
\item But, it is showing the situation for the major language for each country
  \begin{itemize}
  \item Indonesia has over 700 living languages
  \item Australia has over 100 indigenous languages and probably at least that many immigrant languages
  \item Even Czechia recognizes Belarusian, German, Polish, Hungarian, Ukrainian and  Vietnamese as official minority languages
  \end{itemize}
\item We are missing most of the world's languages
  \begin{itemize}
  \item Most wordnets start by exploiting existing resources
    \begin{itemize}
    \item  bootstrap off English
    \end{itemize}
  \item But most of the world's languages do not have lexicons linked
    to English!
  \end{itemize}
\end{itemize}

\end{frame}


% \begin{frame}
% \frametitle{Conclusion of the Introduction}

% \begin{itemize}
%     \item Wordnets are essential resources for understanding the semantic structure of languages.
%     \item They support a wide range of applications in NLP, linguistics, and AI.
%     \item Wordnets are now being developed for many languages, including those with limited resources.
%     \item Multilingual wordnets enable cross-linguistic research and better understanding of language relationships.
%     \item Wordnets contribute to preserving linguistic diversity by documenting and organizing lexical data.
% \end{itemize}
% \end{frame}


 

\section{Low Resource Languages (LRL)}

\begin{frame}
\frametitle{Challenges in Building Wordnets for Low-Resource Languages}

\begin{itemize}
    \item Many languages lack  lexical resources like dictionaries or corpora.
    \item In this case we have to build a foundational lexical database from scratch.
    \item Fieldwork and collaboration with native speakers are essential for data collection.
    \item Cultural and contextual differences make direct translations of concepts difficult.
    \item The lack of standardized orthographies adds complexity to digitization and documentation.
\end{itemize}
\end{frame}

\begin{frame}
\frametitle{Data Collection Challenges}

\begin{itemize}
    \item Workshops, interviews, and manual documentation are often necessary for word collection.
    \item Language consultants may be required for translating and interpreting lexical data.
    \item Typically language experts and computer experts are not the same people.
    \item Word meanings may vary by dialect or region, making consistent data collection difficult.
    \item Language communities with no history of standardisation may disagree as to what should be considered correct
    \item Time and cost constraints limit how much data can be collected.
      \\ also true for well-resourced languages!
\end{itemize}
\end{frame}

\begin{frame}
\frametitle{Semantic Structure Challenges}

\begin{itemize}
    \item Identifying and organizing synsets in under-resourced languages can be difficult.
    \item Concepts may not map directly onto concepts in other languages, such as English.
    \item Cultural concepts and practices often require unique synsets that do not exist in larger languages.
    \item Expert knowledge of the language's lexical semantics is essential for accurate synset creation.
\end{itemize}
\end{frame}

\begin{frame}
\frametitle{Technological Challenges}

\begin{itemize}
    \item Tools and resources for building wordnets, such as NLP software, are not readily available for low-resource languages.
    \item Manual annotation and digitization of handwritten data are time-consuming and error-prone.
    \item Often it is not just that you don't have, for example, a  part-of-speech tagger, but that no-one had yet identified what parts of speech are appropriate
    \item Collaboration between linguists and native speakers can be logistically challenging without reliable power, internet or software tools.
\end{itemize}
\end{frame}

\begin{frame}
\frametitle{Sociolinguistic and Cultural Challenges}

\begin{itemize}
    \item Some communities may resist linguistic documentation efforts due to concerns about language preservation.
    \item Native speakers may prioritize language revitalization over computational resources like wordnets.
    \item The process of documenting a language for a wordnet can introduce external cultural biases.
    \item[★] Collaboration with local communities is essential to ensure cultural respect and accuracy.
    \item[★] Ethical considerations are crucial when engaging with endangered language communities.
\end{itemize}
\end{frame}

\begin{frame}
\frametitle{Challenges in Data Validation and Accuracy}

\begin{itemize}
    \item Data collected from native speakers will contain inconsistencies or errors.
    \item Language experts are often needed to validate synsets and relations.
    \item Regular updates and revisions are necessary as languages evolve or more data is collected.
    \item Data verification is particularly challenging for languages with no formal linguistic documentation, and few or no collections of text.
\end{itemize}
\end{frame}

\begin{frame}
\frametitle{Community Involvement}

\begin{itemize}
    \item Involving the language community is key to successful wordnet development.
    \item Native speakers contribute cultural and contextual insights critical to building accurate wordnets.
    \item Community workshops help ensure that the wordnet reflects the language as spoken by its speakers.
    \item Community members often play a key role in digitizing and validating the wordnet.
    \item Without community involvement, the wordnet may not reflect the real linguistic and cultural landscape.
\end{itemize}
\end{frame}

\begin{frame}{Wordnet is not designed to cover everything}

  \begin{itemize}
  \item The original Princeton wordnet only covered content words: nouns, verbs, adjectives and adverbs in English
    \begin{itemize}
    \item There was no need to cover every word class --- plenty of other dictionaries do that
    \end{itemize}
  \item But if a language has no dictionary, so the wordnet is going to be the \txx{only} lexicon --- then it needs to cover everything the community needs
      \begin{itemize}
      \item Other parts of speech
      \item Usage noted
      \item Audio
      \item Spelling variation
      \end{itemize}
    \item And the data should be as accessible as possible
      \begin{itemize}
      \item to the community
      \item to the field linguists
      \item to other researchers
      \end{itemize}
  \end{itemize}
  
\end{frame}

%\begin{frame}
  
%  \frametitle{Challenges in Maintaining Wordnets}

% \begin{itemize}
%     \item Maintenance requires continuous updates as the language changes or new words are added.
%     \item Low-resource languages may not have sufficient funding or institutional support for long-term projects.
%     \item Finding skilled annotators or linguists for updates can be difficult in smaller communities.
%     \item Technical infrastructure (e.g., servers, hosting) must be sustained over time to keep the wordnet accessible.
%     \item Support for open-source projects like wordnets may diminish without strong community or academic backing.
% \end{itemize}
% \end{frame}

\begin{frame}
\frametitle{Balancing Preservation and Innovation}

\begin{itemize}
    \item Wordnets can support language preservation efforts by documenting lexical data.
    \item At the same time, they promote technological innovation, enabling computational uses of the language.
    \item There is often tension between focusing on language preservation versus creating resources for NLP applications.
    \item Balancing traditional language documentation with new technological tools is a key challenge.
    \item Wordnets offer a way to bridge the gap between preserving linguistic heritage and advancing technology.
\end{itemize}
\end{frame}

%\begin{frame}
% \frametitle{Conclusion of the Challenges}

% \begin{itemize}
%     \item Building wordnets for languages with limited resources involves complex challenges, from data collection to validation.
%     \item Close collaboration with native speakers and language communities is critical for success.
%     \item Despite these challenges, wordnets play a vital role in documenting, preserving, and revitalizing low-resource languages.
%     \item Addressing both linguistic and technical obstacles ensures that these languages remain accessible for future generations.
%     \item Wordnets contribute to the global linguistic diversity and cultural preservation efforts.
% \end{itemize}
% \end{frame}


\section{Case Studies}

\subsection{Abui}

\begin{frame}
\frametitle{Case Study 1: The Abui Wordnet}

\begin{itemize}
    \item Abui (ISO 639-3: abz, abui1241) is a Timor-AlorPantar (TAP) language \citep{Kratochvil:2007}
    \item Spoken by about 17 thousand speakers in an area stretching from the northern to the southern coast in Central Alor
    \item  The work reported here focusses on the variety spoken in the village
of Takalelang at the northern coast.
    \item The Abui wordnet was developed as part of a fieldwork project documenting the language.
    \item Originally the data was stored in SIL toolbox, an NLP tool for field linguists
    \item Then the collection was expanded using Rapid Word Collection (RWC) workshops with local speakers.
    \item Words were collected across domains like plants, animals, emotions, and kinship.
    \item Collected data is in the process of being (manually) digitized and structured into a wordnet format.
\end{itemize}
\end{frame}

\begin{frame}{Alor}
\begin{center}
   \noindent\includegraphics[width=1.0\textwidth]{pics/alor-s.png}
\end{center}
\end{frame}

\begin{frame}
\frametitle{Toolbox to wordnet}
 \begin{tabular}{cc}
  \includegraphics[width=0.4\textwidth]{pics/ToolboxDictFormat-s.png}
   &
 \begin{minipage}{0.5\linewidth}
   \begin{itemize}
   \item Uses the glosses to link to English, Malay and Indonesian wordnets
   \item Intersection in 3 language has an accuracy of 0.99, 2 languages around 0.5 and 1 language 0.35
     \begin{itemize}
     \item Even though Malay and Indonesian are very similar!
     \end{itemize}
   \item Data made available at \url{https://github.com/fanacek/abuiwn}
   \item This wordnet was built using other wordnets
   \end{itemize}
   \hfill \citep{kratochvil-morgado-da-costa-2022-abui}
 \end{minipage}
 \end{tabular}
 \end{frame}


\begin{frame}{Rapid Word Construcion}
  
\begin{itemize}
    \item But we wanted more words!
    \item Three RWC workshops were conducted with over 80 participants across multiple days.
    \item Workshops focused on gathering words within specific semantic domains.
    \item Participants included native speakers from different parts of the community, ensuring dialectal coverage.
    \item Words were first handwritten, then digitized and annotated for semantic relations.
    \item SIL domains were mapped to wordnet concepts, again using translation overlap \citep{Zulhelmy:Bond:Kratochvil:2014,Costa:Kratochvíl:Saad:Delpada:Lanma:Bond:Wolfová:2023}
      \begin{itemize}
      \item because the granularity is very different they do not often match one-to-one
      \item so there is more manual work to be done
      \end{itemize}
    % \item Collaboration with linguists ensured accurate synset creation and semantic alignment.
\end{itemize}
\end{frame}

\begin{frame}{It was a blast}
  \includegraphics[width=1.0\textwidth]{pics/WorkshopSmall-s.jpg}
\end{frame}

\begin{frame}{SIL semantic domain for water (1.3)}
  \begin{tabular}{cc}
    \includegraphics[width=0.4\textwidth]{pics/sil-semdom-water.png}
    &
    \includegraphics[width=0.6\textwidth]{pics/fieldwords-water.png}
      
  \end{tabular}
\end{frame}

\begin{frame}


  \frametitle{Challenges in Building the Abui Wordnet}

\begin{itemize}
\item Data entry was the biggest bottleneck
  \begin{itemize}
  \item Only a native speaker can really digitize the data
  \item There are only a handful who could
  \item They often have other things to do
  \end{itemize}
\item Computational linguist time is also a  bottleneck
  \begin{itemize}
  \item Resource creation is rarely well-funded
  \item It takes second place to other tasks
  \end{itemize}
\item Orthographic variation is also a problem
  \begin{itemize}
  \item The orthography is being refined as we record more data
  \item It is hard to update older data
  \end{itemize}
\end{itemize}
\end{frame}

\begin{frame}
\frametitle{Outcomes of the Abui Wordnet Project}

\begin{itemize}
    \item The Abui wordnet documented over 1,400 synsets and 3,600 senses
    \item A new version with 2,500 synsets is on its way
    \item It serves as a lexical resource for both linguistic research and the local community.
    \item The project supported ongoing language documentation and preservation efforts.
    \item Native speakers have been deeply involved in the project and are training in linguistics
    \item The Abui wordnet was incorporated into the Open Multilingual Wordnet (OMW) project --- you can even find it in hugging face!
\end{itemize}
\end{frame}

\begin{frame}
\frametitle{Significance of the Abui Wordnet}

\begin{itemize}
    \item The wordnet highlights the role of collaborative fieldwork in low-resource language documentation.
    \item It provides a linguistic resource for Abui speakers, researchers, and educators.
    \item The project helps preserve Abui by recording its lexical structure and key cultural concepts.
    \item It represents a model for other Papuan and low-resource languages in similar contexts.
    \item The Abui wordnet demonstrates how modern tools can assist in documenting oral traditions.
    \item We are currently working on annotating a text: \textit{Bukuuting bikaat-bikaat} ``The Speckled Band', which we translated from Indonesian
\\ for this we need a lemmatizer, \ldots
\end{itemize}
\end{frame}

\subsection{Kristang}

\begin{frame}
\frametitle{Case Study 2: The Kristang Wordnet}

\begin{itemize}
    \item Kristang is a critically endangered creole language
    \item Spoken mainly by Portuguese-Eurasian communities in Malacca
and Singapore.
\item There are no more than a few thousand speakers, with more in Malacca than Singapore. 
\item Kristang is a originally derived from Malay and Portuguese
  \begin{itemize}
  \item Vocabulary is largely from Portuguese
  \item Grammar is very close to Malay
  \item Influenced also by Dutch, English and also other Portuguese creoles in Africa, India, South East Asia and China
  \end{itemize}
\end{itemize}
\end{frame}
\begin{frame}{Portuguese Creoles in Asia}
  \includegraphics[width=1.0\textwidth]{pics/portuguese-creoles-in-asia.png}
\end{frame}

\begin{frame}
\frametitle{Language Revitalization and the Kristang Wordnet}

\begin{itemize}
\item The Kristang wordnet is part of a language revitalization effort
  \textit{Kodrah Kristang} ``Awaken Kristang''
\item The goal is to support the
  dwindling speaker base \citep{morgado-da-costa-2020-pinchah}.
    \item The wordnet aims to preserve Kristang's lexicon and support its transmission to younger generations.
    \item The project is led by  local community members, often linguistic students.
    \item Classes and workshops are held to teach Kristang using materials derived from the wordnet.
    \item Workshops are also held to add new words to the wordnet
    \item Community involvement is central, with local speakers helping to curate and validate data, efforts are supported by partnerships with linguistic organizations and local governments.
\end{itemize}
\end{frame}

\begin{frame}{Kodrah Kristang Class}
  \includegraphics[width=1.0\textwidth]{pics/kodrah-kristang.png}
\end{frame}

\begin{frame}
\frametitle{Challenges in Building the Kristang Wordnet}

\begin{itemize}
\item Kristang's limited written tradition required reliance on a wide variety of  language resources.
  \begin{itemize}
  \item paper dictionaries and word lists, including personal collections
  \item linguistic publications with wordlists or glossed text
  \item language documentation work (including  \textit{Kodrah Kristang} and a course at NUS  on Field Methods in Linguistics )
  \item new words and senses produced by \texttt{Jardinggu} ``langarden'', the Kristang lexical incubation project
  \end{itemize}
    \item The small speaker base makes data collection and validation challenging.
   
\end{itemize}
\end{frame}

\begin{frame}
\frametitle{Innovation in the Kristang Wordnet Project}

\begin{itemize}
  \item The project took advantage of existing work in the open multilingual wordnet to also allow an extended inventory of parts of speech:  pronouns, classifiers, exclamatives and so forth \citep{Seah:Bond:2014,daCosta:Bond:2016}
    \item It is also important to make the lexical data available even when not fully integrated
      \begin{itemize}
      \item In order to be useful to the community the Open Kristang
        Wordnet has two layers
      \item The core wordnet with synsets fully linked to the OMW
      \item The extended wordnet that just has bilingual Kristang-English pairs
      \end{itemize}
\end{itemize}
\end{frame}

\begin{frame}
\frametitle{Significance of the Kristang Wordnet}

\begin{itemize}
    \item The wordnet supports efforts to reverse language shift in the Kristang-speaking community.
    \item It highlights the importance of collaboration between linguists and language communities.
    \item The Kristang wordnet can serve as a model for revitalization efforts in other endangered languages.
    \item It demonstrates how lexical resources can play a role in both preservation and education.
    \item The project shows the potential for wordnets in revitalizing endangered creoles and minority languages.
\end{itemize}
\end{frame}

\subsection{Cantonese}

\begin{frame}
\frametitle{Case Study 3: The Cantonese Wordnet}

\begin{itemize}
    \item Cantonese is spoken by millions, but its written tradition is limited compared to Mandarin.
    \item The Cantonese wordnet project aims to provide a lexical resource for this major Chinese variety.
    \item The Cantonese wordnet includes everyday Cantonese vocabulary, colloquialisms, and slang.
    \item It currently has senses, examples, a small sense-tagged corpus \citep{sio-costa-2019-building,sio:morgadodacosta:2022a}
    \item It also has parts-of-speech not in the original wordnet: classifiers and aspect markers
\end{itemize}
\end{frame}
\begin{frame}{Where is Cantonese (and other Chineses) spoken?}
  \includegraphics[height=0.8\textheight]{pics/cantonese-other-chinese-map.png}

  From \citet{Zhenxing:2012}.
\end{frame}


\begin{frame}
\frametitle{Challenges in Building the Cantonese Wordnet}

\begin{itemize}
    \item Cantonese has a strong oral tradition but less standardized written forms.
    \item Because of this there is quite a bit of variation
    \item As Cantonese is a spoken variety, and there is some
      variation in pronunciation, we decided to add audio to the wordnet
      \begin{itemize}
      \item This takes advantage of the extensions to the wordnet format from 2020
      \item As far as we know we are the first people to add sound
      \item We recorded some data ourselves
      \item We harvested some data from Wikicommons
        \begin{itemize}
        \item Many entries were Mandarin with Cantonese pronunciation, not Cantonese
        \item It is difficult to distinguish them
        \item But essential to do so
        \end{itemize}
      \end{itemize}
    \item It was sometimes hard to distinguish words from phrases
      \begin{itemize}
      \item Written Chinese does not show word boundaries
      \item Linguists disagree on what should be a word
      \end{itemize}

\end{itemize}
\end{frame}


\begin{frame}{The OMW interface showing audio}
    \includegraphics[width=1.0\textwidth]{pics/nei5.png}

    \begin{itemize}
    \item This shows a word with two pronunciations
      \begin{itemize}
      \item nei5 is the standard pronunciation
      \item lei5 is the `lazy' pronunciation \citep{Chen:2018}
      \end{itemize}
    \end{itemize}
    
\end{frame}


\begin{frame}
\frametitle{Outcomes of the Cantonese Wordnet Project}

\begin{itemize}
\item The Cantonese wordnet is a carefully curated resource 
  \begin{itemize}
  \item 6,200 concepts
  \item 17,350 senses
  \item 2,138 audio examples, covering 2,859 senses
    % (+ 5250 964) 6214 (+ 16300 1050) 17350
  \end{itemize}
    \item The project aids in preserving Cantonese as a separate linguistic entity from Mandarin.
    \item The wordnet has been incorporated into multilingual wordnets, supporting cross-linguistic research.
    \item It serves as a basis for future research on Cantonese language technology development.
\end{itemize}
\end{frame}

\begin{frame}
\frametitle{Significance of the Cantonese Wordnet}

\begin{itemize}
    \item The Cantonese wordnet serves as a reference for the vocabulary of the language.
    \item It supports the development of linguistic resources for non-Mandarin Chinese varieties.
    \item The wordnet helps capture the cultural and linguistic identity of Cantonese speakers.
    \item It contributes to the broader effort of documenting and preserving regional Chinese languages.
    \item The project highlights the importance of creating linguistic resources for spoken varieties.
\end{itemize}
\end{frame}

\section{Advantages and disadvantages of using wordnets for LRL}

\begin{frame}
\frametitle{Advantages of Wordnets for Low-Resource Languages}

\begin{itemize}
    \item Wordnets help document and preserve linguistic data for endangered or under-documented languages.
    \item They enable cross-linguistic comparison by aligning synsets with wordnets of other languages.
    \item Wordnets provide a structured, searchable resource that benefits researchers and language communities.
    \item They can be used to support language revitalization efforts and to create educational materials.
      \begin{itemize}
      \item Having an online presence can boost the social-status of a language
      \end{itemize}
    \item Wordnets help build other wordnets
      \begin{itemize}
      \item More languages allows better bootstrapping
      \item Phenomena covered in one language make it easier for the next
      \item The more data there is, the better the descriptions become
      \end{itemize}
\end{itemize}
\end{frame}

\begin{frame}
\frametitle{Cultural and Linguistic Preservation}

\begin{itemize}
\item Wordnets should record and organize culturally significant concepts that may not exist in other languages.
  \begin{itemize}
  \item This has not been done as much as it should
  \item Adding and describing new concepts is the next great challenge!
  \end{itemize}
    \item They create a permanent record of a language's lexicon, supporting long-term preservation.
    \item Wordnets reflect the unique cultural and cognitive world of speakers, documenting traditional knowledge.
    \item In endangered languages, wordnets can capture the lexical heritage before it disappears.
    \item They facilitate the transmission of traditional vocabulary to younger generations.
\end{itemize}
\end{frame}

\begin{frame}
\frametitle{Technological Benefits}

\begin{itemize}
    \item Wordnets support the development of language technology for low-resource languages.
    \item They provide a foundation for tools like speech recognition, machine translation, and sentiment analysis.
    \item Wordnets enable the development of intelligent systems that can understand semantic relationships in texts.
    \item They help create more inclusive digital resources, bringing under-represented languages online.
    \item Open-source platforms like the Open Multilingual Wordnet (OMW) facilitate easy integration into existing projects.
\end{itemize}
\end{frame}

\begin{frame}
\frametitle{Educational and Revitalization Efforts}

\begin{itemize}
    \item Wordnets can be used to create language learning materials and dictionaries for education.
    \item They provide an organized resource for teaching both linguistic structure and vocabulary.
    \item Revitalization programs benefit from wordnets by using them to foster literacy and language proficiency.
    \item Wordnets allow the community to access and engage with their language in new technological formats.
    \item They offer a digital resource for the continued learning and use of the language in modern contexts.
\end{itemize}
\end{frame}

\begin{frame}
\frametitle{Advantages in Linguistic Research}

\begin{itemize}
    \item Wordnets allow linguists to analyze and compare the lexical structure of different languages.
    \item They provide insights into how languages encode meanings and semantic relationships.
    \item Multilingual wordnets facilitate cross-linguistic research on polysemy, metaphor, and lexical relations.
    \item Wordnets serve as an important resource in typological studies and language contact research.
      \begin{itemize}
      \item For example we used the Moroccan Arabic wordnet to investigate shared vocabulary with Standard Arabic, French and Italian \citep{Mrini:Bond:2018}
      \end{itemize}
    \item They enable the study of cognitive aspects of language, such as categorization and conceptualization.
\end{itemize}
\end{frame}

\section{Thanks}

\begin{frame}{Thanks and disclaimer}

  \begin{itemize}
  \item Thank you to John McCrae and the Teanga Project for inviting me to give this talk
  \item Thanks to all the many people who have worked on these
    resources, especially Luís Morgado da Costa, František Kratochvíl
    and Joanna Ut-Seong Sio.
  \item I have talked about projects my lab has been involved in, as I know them best, but there are many other wordnets for LRL
    \begin{itemize}
    \item Amharic, Kurdish, Mansi, Moroccan Arabic, Sardinian, ASL, Uzbek, Welsh, \ldots
    \end{itemize}
  \item There are also wordnets for ancient languages \begin{itemize}
    \item Ancient Greek, Coptic, Latin, Qin Chinese, Sanskrit\ldots
    \end{itemize}
  \end{itemize}


  
\end{frame}


\begin{frame}[allowframebreaks]
        \frametitle{References}
        \bibliographystyle{aclnat}
        \bibliography{abb,mtg,ling,nlp}
\end{frame}



\end{document}

%%% Local Variables: 
%%% coding: utf-8
%%% mode: latex
%%% TeX-PDF-mode: t
%%% TeX-engine: xetex
%%% End: 
